\hypertarget{standard-template-library-core}{%
\chapter{STANDARD TEMPLATE LIBRARY
CORE}\label{standard-template-library-core}}

\hypertarget{the-standard-template-library-stl-core}{%
\chapter{THE STANDARD TEMPLATE LIBRARY (STL)
CORE}\label{the-standard-template-library-stl-core}}

\hypertarget{c9803-standard-library}{%
\chapter{C++98/03 STANDARD LIBRARY}\label{c9803-standard-library}}

\hypertarget{standard-template-library-stl}{%
\section{Standard Template Library
(STL)}\label{standard-template-library-stl}}

\begin{center}\rule{0.5\linewidth}{0.5pt}\end{center}

\hypertarget{introduction-to-stl}{%
\section{Introduction to STL}\label{introduction-to-stl}}

The Standard Template Library (STL) is a collection of template classes
and functions that provide: - \textbf{Containers} - Data structures to
hold objects (e.g., vectors, lists, maps). - \textbf{Iterators} -
Objects to traverse containers (generalization of pointers). -
\textbf{Algorithms} - Functions to manipulate data (e.g., sorting,
searching). - \textbf{Function Objects (Functors)} - Objects that act
like functions.

\hypertarget{key-advantages}{%
\subsection{Key Advantages}\label{key-advantages}}

\begin{itemize}
\tightlist
\item
  \textbf{Generic Programming}: Write code that works with any data
  type.
\item
  \textbf{Performance}: Heavily optimized implementations.
\item
  \textbf{Reusability}: Standardized components prevent reinventing the
  wheel.
\item
  \textbf{Type Safety}: Templates ensure type correctness at compile
  time.
\end{itemize}

\begin{center}\rule{0.5\linewidth}{0.5pt}\end{center}

\hypertarget{stl-components-overview}{%
\section{STL Components Overview}\label{stl-components-overview}}

\begin{lstlisting}
          STL (Standard Template Library)  
                       |
   -----------------------------------------
   |              |            |           |
CONTAINERS     ITERATORS    ALGORITHMS   FUNCTORS
   |              |            |           |
Sequence       Input        Searching    Predicates
Associative    Output       Sorting      Comparators
Adapters       Forward      Modifying      
               Bidir        Numeric      
               Random       
\end{lstlisting}

\begin{center}\rule{0.5\linewidth}{0.5pt}\end{center}

\hypertarget{professional-notes-stl-core-depth}{%
\subsection{Professional Notes: STL Core
Depth}\label{professional-notes-stl-core-depth}}

\hypertarget{iterators-the-bridge-between-algorithms-and-containers}{%
\subsubsection{1. Iterators: The Bridge between Algorithms and
Containers}\label{iterators-the-bridge-between-algorithms-and-containers}}

Iterators provide a uniform interface for traversing data. *
\textbf{Input/Output}: Single pass, read or write once. *
\textbf{Forward}: Multiple passes, read/write, move forward only (e.g.,
\passthrough{\lstinline!std::forward\_list!}). * \textbf{Bidirectional}:
Move forward and backward (e.g., \passthrough{\lstinline!std::list!},
\passthrough{\lstinline!std::map!}). * \textbf{Random Access}: Jump to
any element in constant time (e.g.,
\passthrough{\lstinline!std::vector!},
\passthrough{\lstinline!std::deque!}).

\textbf{Godhood Tip}: Prefer prefix increment
(\passthrough{\lstinline!++it!}) over postfix increment
(\passthrough{\lstinline!it++!}) for iterators. Postfix creates a
temporary copy of the iterator, which can be costly for complex types.

\hypertarget{container-choice-and-memory-layout}{%
\subsubsection{2. Container Choice and Memory
Layout}\label{container-choice-and-memory-layout}}

\begin{itemize}
\tightlist
\item
  \textbf{\passthrough{\lstinline!std::vector!}}: The gold standard.
  Contiguous memory ensures high \textbf{Cache Locality}. Always
  \passthrough{\lstinline!reserve()!} if you know the final size to
  avoid reallocations.
\item
  \textbf{\passthrough{\lstinline!std::deque!}}: A ``Double-Ended
  Queue.'' Implemented as a sequence of fixed-size memory blocks. Offers
  \(O(1)\) at both ends but is slower for random access than vector.
\item
  \textbf{\passthrough{\lstinline!std::list!}}: Doubly linked list.
  \(O(1)\) insertions anywhere if you have the iterator, but terrible
  cache locality and high memory overhead per element (2 pointers).
\end{itemize}

\hypertarget{associative-containers-and-custom-comparators}{%
\subsubsection{3. Associative Containers and Custom
Comparators}\label{associative-containers-and-custom-comparators}}

\passthrough{\lstinline!std::map!} and
\passthrough{\lstinline!std::set!} are typically implemented as
\textbf{Red-Black Trees}. * \textbf{Complexity}: \(O(\log N)\) for all
major operations. * \textbf{Comparators}: You can provide a custom
function or functor to define the ordering.

\begin{lstlisting}[language={C++}]
struct CaseInsensitiveCompare {
    bool operator()(const std::string& a, const std::string& b) const {
        return strcasecmp(a.c_str(), b.c_str()) < 0;
    }
};
std::set<std::string, CaseInsensitiveCompare> my_set;
\end{lstlisting}

\begin{longtable}[]{@{}
  >{\raggedright\arraybackslash}p{(\columnwidth - 0\tabcolsep) * \real{0.06}}@{}}
\toprule
\endhead
\#\#\# Professional Notes: Data Structures Internals \\
\#\#\#\# 1. Binary Search Trees (std::map, std::set) Typically
implemented as \textbf{Red-Black Trees}. * \textbf{Self-Balancing}:
Ensures \(O(\log N)\) height. * \textbf{Node Overhead}: Each element is
stored in a separate node with pointers to parent and children, plus a
color bit. \\
\#\#\#\# 2. Hash Tables (std::unordered\_map) Typically implemented as
an array of buckets (linked lists). * \textbf{Load Factor}: When the
number of elements exceeds
\passthrough{\lstinline!bucket\_count * max\_load\_factor!}, the table
is rehashed (size doubled). * \textbf{Hash Collisions}: Handled via
chaining (linked lists) or open addressing. \\
\#\#\#\# 3. Heap (std::priority\_queue) Implemented using
\passthrough{\lstinline!std::make\_heap!},
\passthrough{\lstinline!std::push\_heap!}, and
\passthrough{\lstinline!std::pop\_heap!} on an underlying
\passthrough{\lstinline!std::vector!}. * \textbf{Invariant}: The parent
is always greater than (or equal to) its children. \\
\bottomrule
\end{longtable}

\begin{center}\rule{0.5\linewidth}{0.5pt}\end{center}

\hypertarget{containers---complete-reference}{%
\section{CONTAINERS - COMPLETE
REFERENCE}\label{containers---complete-reference}}

\hypertarget{container-characteristics-c98}{%
\subsection{Container Characteristics
(C++98)}\label{container-characteristics-c98}}

\begin{longtable}[]{@{}lllllll@{}}
\toprule
Container & Type & Insert & Delete & Search & Random Access & Memory \\
\midrule
\endhead
\passthrough{\lstinline!vector!} & Sequence & O(n) & O(n) & O(n) & O(1)
& Contiguous \\
\passthrough{\lstinline!list!} & Sequence & O(1) & O(1) & O(n) & - &
Scattered \\
\passthrough{\lstinline!deque!} & Sequence & O(n) & O(n) & O(n) & O(1) &
Chunks \\
\passthrough{\lstinline!map!} & Associative & O(log n) & O(log n) &
O(log n) & - & Tree \\
\passthrough{\lstinline!set!} & Associative & O(log n) & O(log n) &
O(log n) & - & Tree \\
\passthrough{\lstinline!multimap!} & Associative & O(log n) & O(log n) &
O(log n) & - & Tree \\
\passthrough{\lstinline!multiset!} & Associative & O(log n) & O(log n) &
O(log n) & - & Tree \\
\passthrough{\lstinline!priority\_queue!} & Adapter & O(log n) & O(log
n) & - & - & Heap \\
\passthrough{\lstinline!queue!} & Adapter & O(1) & O(1) & - & - & - \\
\passthrough{\lstinline!stack!} & Adapter & O(1) & O(1) & - & - & - \\
\bottomrule
\end{longtable}

\begin{center}\rule{0.5\linewidth}{0.5pt}\end{center}

\hypertarget{vector---dynamic-array}{%
\section{1.1 VECTOR - Dynamic Array}\label{vector---dynamic-array}}

\hypertarget{what-is-vector}{%
\subsection{What is Vector?}\label{what-is-vector}}

A dynamic array that grows automatically. Use this as your default
container.

\hypertarget{declaration-initialization}{%
\subsection{Declaration \&
Initialization}\label{declaration-initialization}}

\begin{lstlisting}[language={C++}]
#include <vector>
using namespace std;

// Empty vector
vector<int> v1;

// Vector with initial size (10 elements initialized to 0)
vector<int> v2(10);

// Vector with initial size and value (5 elements, all 10)
vector<int> v3(5, 10);

// Copy constructor
vector<int> v4(v3);

// From array (using pointers)
int arr[] = {1, 2, 3, 4, 5};
vector<int> v5(arr, arr + 5); 
\end{lstlisting}

\hypertarget{accessing-elements}{%
\subsection{Accessing Elements}\label{accessing-elements}}

\begin{lstlisting}[language={C++}]
vector<int> v;
v.push_back(10);
v.push_back(20);

// 1. Using operator[] (No bounds check)
cout << v[0] << "\n";  // 10

// 2. Using at() (With bounds check, throws std::out_of_range)
cout << v.at(1) << "\n"; // 20

// 3. Front and back
cout << v.front() << "\n"; // 10
cout << v.back() << "\n";  // 20
\end{lstlisting}

\hypertarget{modifying-elements}{%
\subsection{Modifying Elements}\label{modifying-elements}}

\begin{lstlisting}[language={C++}]
vector<int> v;

// Adding elements
v.push_back(10);
v.push_back(20);
v.push_back(30); // {10, 20, 30}

// Inserting elements (Iterators required)
v.insert(v.begin() + 1, 15); // {10, 15, 20, 30}

// Removing elements
v.pop_back();           // Removes 30
v.erase(v.begin() + 1); // Removes 15

// Clearing
v.clear();              // Empty
\end{lstlisting}

\hypertarget{size-capacity}{%
\subsection{Size \& Capacity}\label{size-capacity}}

\begin{lstlisting}[language={C++}]
vector<int> v(10);

cout << v.size() << "\n";      // 10
cout << v.capacity() << "\n";  // >= 10
cout << v.empty() << "\n";     // 0 (false)

// Reserve space (Optimization to avoid reallocations)
v.reserve(100); 

// Resize (Changes size, fills new elements with default or value)
v.resize(20);    // Size becomes 20
v.resize(5);     // Size becomes 5, extra elements destroyed
\end{lstlisting}

\hypertarget{iterating-through-vector}{%
\subsection{Iterating Through Vector}\label{iterating-through-vector}}

\begin{lstlisting}[language={C++}]
vector<int> v;
v.push_back(10); v.push_back(20); v.push_back(30);

// 1. Index loop
for (size_t i = 0; i < v.size(); ++i) {
    cout << v[i] << " ";
}

// 2. Iterator loop (Recommended)
for (vector<int>::iterator it = v.begin(); it != v.end(); ++it) {
    cout << *it << " ";
}

// 3. Reverse Iterator
for (vector<int>::reverse_iterator rit = v.rbegin(); rit != v.rend(); ++rit) {
    cout << *rit << " ";
}
\end{lstlisting}

\begin{center}\rule{0.5\linewidth}{0.5pt}\end{center}

\hypertarget{deque---double-ended-queue}{%
\section{1.2 DEQUE - Double Ended
Queue}\label{deque---double-ended-queue}}

\hypertarget{what-is-deque}{%
\subsection{What is Deque?}\label{what-is-deque}}

Fast insertion/deletion at both ends. Unlike vector, storage is not
guaranteed to be contiguous.

\begin{lstlisting}[language={C++}]
#include <deque>
using namespace std;

deque<int> dq;

dq.push_back(10);
dq.push_front(5);  // {5, 10}

dq.pop_back();     // {5}
dq.pop_front();    // {}

// Supports random access
dq.push_back(100);
cout << dq[0] << "\n";
\end{lstlisting}

\begin{center}\rule{0.5\linewidth}{0.5pt}\end{center}

\hypertarget{list---doubly-linked-list}{%
\section{1.3 LIST - Doubly Linked
List}\label{list---doubly-linked-list}}

\hypertarget{what-is-list}{%
\subsection{What is List?}\label{what-is-list}}

Efficient insertion/deletion anywhere (O(1) if iterator is known). No
random access.

\begin{lstlisting}[language={C++}]
#include <list>
using namespace std;

list<int> lst;
lst.push_back(10);
lst.push_front(5);

// Insert at position
list<int>::iterator it = lst.begin();
++it; // Move to second position
lst.insert(it, 7); // {5, 7, 10}

// Remove
lst.remove(7); // Remove all elements with value 7

// Unique (Removes consecutive duplicates)
lst.push_back(10);
lst.unique(); // {5, 10} (second 10 removed)

// Sort
lst.sort(); // Internal sort (std::sort doesn't work on lists)
\end{lstlisting}

\begin{center}\rule{0.5\linewidth}{0.5pt}\end{center}

\hypertarget{map---sorted-key-value-pairs}{%
\section{1.4 MAP - Sorted Key-Value
Pairs}\label{map---sorted-key-value-pairs}}

\hypertarget{what-is-map}{%
\subsection{What is Map?}\label{what-is-map}}

Associative container storing key-value pairs sorted by key. Implemented
as a Balanced BST (usually Red-Black Tree).

\begin{lstlisting}[language={C++}]
#include <map>
#include <string>
using namespace std;

map<string, int> ages;

// Insertion
ages["Alice"] = 30;
ages["Bob"] = 25;
ages.insert(make_pair("Charlie", 28));

// Access / Search
cout << ages["Alice"] << "\n"; // 30

map<string, int>::iterator it = ages.find("Bob");
if (it != ages.end()) {
    cout << "Bob is " << it->second << " years old.\n";
}

// Iteration (Sorted by key)
for (it = ages.begin(); it != ages.end(); ++it) {
    cout << it->first << ": " << it->second << "\n";
}
\end{lstlisting}

\begin{center}\rule{0.5\linewidth}{0.5pt}\end{center}

\hypertarget{set---sorted-unique-keys}{%
\section{1.5 SET - Sorted Unique Keys}\label{set---sorted-unique-keys}}

\hypertarget{what-is-set}{%
\subsection{What is Set?}\label{what-is-set}}

Stores unique elements sorted by value.

\begin{lstlisting}[language={C++}]
#include <set>
using namespace std;

set<int> s;
s.insert(10);
s.insert(5);
s.insert(10); // Duplicate ignored

// {5, 10}

if (s.find(5) != s.end()) {
    cout << "5 is present.\n";
}

// Range erase
s.erase(s.begin()); // Removes 5
\end{lstlisting}

\begin{center}\rule{0.5\linewidth}{0.5pt}\end{center}

\hypertarget{multimap-multiset}{%
\section{1.6 MULTIMAP \& MULTISET}\label{multimap-multiset}}

Allow duplicate keys (multimap) or duplicate values (multiset).

\begin{lstlisting}[language={C++}]
#include <map>
using namespace std;

multimap<string, int> scores;
scores.insert(make_pair("Alice", 90));
scores.insert(make_pair("Alice", 85)); // Allowed

// Finding all values for a key
pair<multimap<string, int>::iterator, multimap<string, int>::iterator> range;
range = scores.equal_range("Alice");

for (multimap<string, int>::iterator it = range.first; it != range.second; ++it) {
    cout << it->second << " ";
}
// Output: 90 85
\end{lstlisting}

\begin{center}\rule{0.5\linewidth}{0.5pt}\end{center}

\hypertarget{stack-adapter}{%
\section{1.7 STACK (Adapter)}\label{stack-adapter}}

LIFO (Last In, First Out). Wraps \passthrough{\lstinline!deque!}
(default), \passthrough{\lstinline!vector!}, or
\passthrough{\lstinline!list!}.

\begin{lstlisting}[language={C++}]
#include <stack>
using namespace std;

stack<int> st;
st.push(10);
st.push(20);

cout << st.top() << "\n"; // 20
st.pop();                 // Removes 20
\end{lstlisting}

\begin{center}\rule{0.5\linewidth}{0.5pt}\end{center}

\hypertarget{queue-adapter}{%
\section{1.8 QUEUE (Adapter)}\label{queue-adapter}}

FIFO (First In, First Out). Wraps \passthrough{\lstinline!deque!}
(default) or \passthrough{\lstinline!list!}.

\begin{lstlisting}[language={C++}]
#include <queue>
using namespace std;

queue<int> q;
q.push(10);
q.push(20);

cout << q.front() << "\n"; // 10
q.pop();                   // Removes 10
\end{lstlisting}

\begin{center}\rule{0.5\linewidth}{0.5pt}\end{center}

\hypertarget{priority_queue-adapter}{%
\section{1.9 PRIORITY\_QUEUE (Adapter)}\label{priority_queue-adapter}}

Sorted queue (Heap). Max element is always at
\passthrough{\lstinline!top()!}.

\begin{lstlisting}[language={C++}]
#include <queue>
#include <vector>
#include <functional> // for greater<int>
using namespace std;

// Max Heap (Default)
priority_queue<int> pq;
pq.push(10);
pq.push(30);
pq.push(20);

cout << pq.top() << "\n"; // 30

// Min Heap
priority_queue<int, vector<int>, greater<int> > min_pq;
min_pq.push(10);
min_pq.push(30);

cout << min_pq.top() << "\n"; // 10
\end{lstlisting}

\begin{center}\rule{0.5\linewidth}{0.5pt}\end{center}

\hypertarget{iterators}{%
\section{ITERATORS}\label{iterators}}

Iterators are the glue between Containers and Algorithms.

\hypertarget{categories}{%
\subsection{Categories}\label{categories}}

\begin{enumerate}
\def\labelenumi{\arabic{enumi}.}
\tightlist
\item
  \textbf{Input}: Read-only, single pass (e.g.,
  \passthrough{\lstinline!istream\_iterator!}).
\item
  \textbf{Output}: Write-only, single pass (e.g.,
  \passthrough{\lstinline!ostream\_iterator!}).
\item
  \textbf{Forward}: Read/Write, multi-pass (e.g.,
  \passthrough{\lstinline!slist!} - non-std).
\item
  \textbf{Bidirectional}: Forward + Backward (e.g.,
  \passthrough{\lstinline!list!}, \passthrough{\lstinline!map!},
  \passthrough{\lstinline!set!}).
\item
  \textbf{Random Access}: O(1) jump (e.g.,
  \passthrough{\lstinline!vector!}, \passthrough{\lstinline!deque!},
  arrays).
\end{enumerate}

\hypertarget{operations}{%
\subsection{Operations}\label{operations}}

\begin{itemize}
\tightlist
\item
  \passthrough{\lstinline!*it!}: Dereference.
\item
  \passthrough{\lstinline!++it!}: Advance.
\item
  \passthrough{\lstinline!--it!}: Retreat (Bidirectional+).
\item
  \passthrough{\lstinline!it + n!}: Jump (Random Access only).
\item
  \passthrough{\lstinline!it1 - it2!}: Distance (Random Access only).
\end{itemize}

\begin{lstlisting}[language={C++}]
vector<int> v(5, 1); 
vector<int>::iterator it = v.begin();
advance(it, 2); // Move 2 steps
cout << distance(v.begin(), it); // 2
\end{lstlisting}

\begin{center}\rule{0.5\linewidth}{0.5pt}\end{center}

\hypertarget{algorithms}{%
\section{ALGORITHMS}\label{algorithms}}

Found in \passthrough{\lstinline!<algorithm>!} and
\passthrough{\lstinline!<numeric>!}. They work on iterator ranges
\passthrough{\lstinline![first, last)!}.

\hypertarget{non-modifying}{%
\subsection{1. Non-Modifying}\label{non-modifying}}

\begin{itemize}
\tightlist
\item
  \passthrough{\lstinline!find(begin, end, val)!}
\item
  \passthrough{\lstinline!count(begin, end, val)!}
\item
  \passthrough{\lstinline!equal(b1, e1, b2)!}
\item
  \passthrough{\lstinline!search(b1, e1, b2, e2)!}
\end{itemize}

\begin{center}\rule{0.5\linewidth}{0.5pt}\end{center}

\hypertarget{professional-notes-algorithm-mastery}{%
\subsection{Professional Notes: Algorithm
Mastery}\label{professional-notes-algorithm-mastery}}

\hypertarget{range-integrity-and-end-iterators}{%
\subsubsection{1. Range Integrity and End
Iterators}\label{range-integrity-and-end-iterators}}

All STL algorithms operate on half-open ranges
\passthrough{\lstinline![first, last)!}. * \textbf{The ``Last''
Iterator}: Points \emph{beyond} the last element. Dereferencing it is
\textbf{Undefined Behavior}. * \textbf{Empty Range}: If
\passthrough{\lstinline!first == last!}, the range is empty. Algorithms
correctly handle this (e.g., \passthrough{\lstinline!find!} returns
\passthrough{\lstinline!last!}).

\hypertarget{sorting-and-stability}{%
\subsubsection{2. Sorting and Stability}\label{sorting-and-stability}}

\begin{itemize}
\tightlist
\item
  \textbf{\passthrough{\lstinline!std::sort!}}: Usually implemented as
  \textbf{Introsort} (Hybrid of Quicksort, Heapsort, and Insertion
  Sort). Average complexity \(O(N \log N)\).
\item
  \textbf{\passthrough{\lstinline!std::stable\_sort!}}: Maintains the
  relative order of equal elements. Requires extra memory for its work
  buffer.
\item
  \textbf{\passthrough{\lstinline!std::partial\_sort!}}: Find the top
  \(K\) elements without sorting the whole range.
\end{itemize}

\hypertarget{the-lambda-evolution-c11}{%
\subsubsection{3. The Lambda Evolution
(C++11)}\label{the-lambda-evolution-c11}}

Algorithms are most powerful when combined with lambdas:

\begin{lstlisting}[language={C++}]
// Find first even number
auto it = std::find_if(vec.begin(), vec.end(), [](int x){ return x % 2 == 0; });
\end{lstlisting}

\hypertarget{binary-search-and-sorted-ranges}{%
\subsubsection{4. Binary Search and Sorted
Ranges}\label{binary-search-and-sorted-ranges}}

Functions like \passthrough{\lstinline!binary\_search!},
\passthrough{\lstinline!lower\_bound!}, and
\passthrough{\lstinline!upper\_bound!} require the range to be
\textbf{sorted} or at least partitioned by the search value. Using them
on unsorted ranges is UB.

\begin{center}\rule{0.5\linewidth}{0.5pt}\end{center}

\hypertarget{modifying}{%
\subsection{2. Modifying}\label{modifying}}

\begin{itemize}
\tightlist
\item
  \passthrough{\lstinline!copy(b1, e1, b2)!}
\item
  \passthrough{\lstinline!transform(b1, e1, out, op)!}
\item
  \passthrough{\lstinline!replace(b1, e1, old, new)!}
\item
  \passthrough{\lstinline!fill(b, e, val)!}
\item
  \passthrough{\lstinline!swap(a, b)!}
\item
  \passthrough{\lstinline!reverse(b, e)!}
\item
  \passthrough{\lstinline!rotate(b, mid, e)!}
\item
  \passthrough{\lstinline!random\_shuffle(b, e)!}
\end{itemize}

\hypertarget{sorting}{%
\subsection{3. Sorting}\label{sorting}}

\begin{itemize}
\tightlist
\item
  \passthrough{\lstinline!sort(b, e)!}: O(N log N).
\item
  \passthrough{\lstinline!stable\_sort(b, e)!}: Preserves order of equal
  elements.
\item
  \passthrough{\lstinline!partial\_sort(b, mid, e)!}: Top K elements.
\end{itemize}

\hypertarget{binary-search-on-sorted-ranges}{%
\subsection{4. Binary Search (On Sorted
Ranges)}\label{binary-search-on-sorted-ranges}}

\begin{itemize}
\tightlist
\item
  \passthrough{\lstinline!binary\_search(b, e, val)!}: Returns bool.
\item
  \passthrough{\lstinline!lower\_bound(b, e, val)!}: First element
  \textgreater= val.
\item
  \passthrough{\lstinline!upper\_bound(b, e, val)!}: First element
  \textgreater{} val.
\end{itemize}

\hypertarget{set-operations-on-sorted-ranges}{%
\subsection{5. Set Operations (On Sorted
Ranges)}\label{set-operations-on-sorted-ranges}}

\begin{itemize}
\tightlist
\item
  \passthrough{\lstinline!set\_union!},
  \passthrough{\lstinline!set\_intersection!},
  \passthrough{\lstinline!set\_difference!}.
\end{itemize}

\hypertarget{numeric}{%
\subsection{6. Numeric ()}\label{numeric}}

\begin{itemize}
\tightlist
\item
  \passthrough{\lstinline!accumulate(b, e, init)!}: Sum.
\item
  \passthrough{\lstinline!inner\_product!}: Dot product.
\item
  \passthrough{\lstinline!adjacent\_difference!}.
\end{itemize}

\hypertarget{example-sort-and-find}{%
\subsection{Example: Sort and Find}\label{example-sort-and-find}}

\begin{lstlisting}[language={C++}]
#include <algorithm>
#include <vector>
#include <iostream>
using namespace std;

bool descending(int a, int b) {
    return a > b;
}

int main() {
    int arr[] = {5, 2, 9, 1, 5, 6};
    vector<int> v(arr, arr + 6);

    // Sort ascending
    sort(v.begin(), v.end());

    // Binary Search
    if (binary_search(v.begin(), v.end(), 9)) {
        cout << "Found 9\n";
    }

    // Sort descending with predicate
    sort(v.begin(), v.end(), descending);

    return 0;
}
\end{lstlisting}

\begin{center}\rule{0.5\linewidth}{0.5pt}\end{center}

\hypertarget{strings-stdstring}{%
\section{STRINGS (std::string)}\label{strings-stdstring}}

A specialization of \passthrough{\lstinline!basic\_string<char>!}.
Replaces C-style \passthrough{\lstinline!char*!} strings.

\begin{lstlisting}[language={C++}]
#include <string>
#include <iostream>
using namespace std;

string s = "Hello";
s += " World"; // Concatenation

// Substring
string sub = s.substr(0, 5); // "Hello"

// Find
size_t pos = s.find("World");
if (pos != string::npos) {
    cout << "Found at " << pos << "\n";
}

// C-string compatibility
const char* cstr = s.c_str();
\end{lstlisting}

\begin{center}\rule{0.5\linewidth}{0.5pt}\end{center}

\hypertarget{streams-iostream}{%
\section{STREAMS (IOSTREAM)}\label{streams-iostream}}

\hypertarget{file-io}{%
\subsection{File I/O ()}\label{file-io}}

\begin{lstlisting}[language={C++}]
#include <fstream>
#include <iostream>
using namespace std;

int main() {
    // Write
    ofstream out("test.txt");
    if (out) {
        out << "Line 1" << endl;
        out << 123 << endl;
    }
    out.close();

    // Read
    ifstream in("test.txt");
    string line;
    int num;
    
    if (in >> line >> num) { // Reads "Line" then fails on "1" vs int? No. 
                             // "Line" goes to line. "1" goes to num? 
                             // Need careful parsing.
    }
    // Better: getline
    getline(in, line);
    
    return 0;
}
\end{lstlisting}

\hypertarget{string-streams}{%
\subsection{String Streams ()}\label{string-streams}}

Useful for parsing strings or formatting.

\begin{lstlisting}[language={C++}]
#include <sstream>
using namespace std;

// Int to String
int x = 42;
stringstream ss;
ss << x;
string s = ss.str();

// String to Int
string s2 = "100";
stringstream ss2(s2);
int y;
ss2 >> y;
\end{lstlisting}

\begin{center}\rule{0.5\linewidth}{0.5pt}\end{center}

\hypertarget{functors-function-objects}{%
\section{FUNCTORS (Function Objects)}\label{functors-function-objects}}

Classes that overload \passthrough{\lstinline!operator()!}. Used by
algorithms.

\begin{lstlisting}[language={C++}]
struct AddX {
    int x;
    AddX(int val) : x(val) {}
    
    int operator()(int y) const {
        return x + y;
    }
};

// Usage with transform
vector<int> v(5, 1); // {1, 1, 1, 1, 1}
transform(v.begin(), v.end(), v.begin(), AddX(10));
// v is now {11, 11, 11, 11, 11}
\end{lstlisting}

This covers the foundational C++98 Standard Library components necessary
for mastery.

\hypertarget{advanced-strings}{%
\chapter{ADVANCED STRINGS}\label{advanced-strings}}

\hypertarget{string-manipulation}{%
\section{5.1 String Manipulation}\label{string-manipulation}}

\begin{lstlisting}[language={C++}]
#include <iostream>
#include <string>
#include <cstring>
using namespace std;

int main() {
    string s = "Hello World";
    
    // Length and capacity
    cout << "Length: " << s.length() << endl;
    cout << "Capacity: " << s.capacity() << endl;
    
    // Access characters
    cout << "First char: " << s[0] << endl;
    cout << "Last char: " << s[s.length() - 1] << endl;
    
    // Finding substrings
    size_t pos = s.find("World");
    if (pos != string::npos) {
        cout << "Found at position: " << pos << endl;
    }
    
    // Replace
    s.replace(6, 5, "C++");
    cout << s << endl;  // Hello C++
    
    // Insert
    s.insert(5, " there");
    cout << s << endl;  // Hello there C++
    
    // Erase
    s.erase(5, 6);
    cout << s << endl;  // Hello C++
    
    // Substring
    cout << s.substr(0, 5) << endl;  // Hello
    
    // Reverse
    reverse(s.begin(), s.end());
    cout << s << endl;  // ++C olleH
    
    return 0;
}
\end{lstlisting}

\hypertarget{string-conversion}{%
\section{5.2 String Conversion}\label{string-conversion}}

\begin{lstlisting}[language={C++}]
#include <iostream>
#include <string>
#include <sstream>
#include <cstdlib>
using namespace std;

int main() {
    // String to number (C style)
    string s1 = "42";
    int num = atoi(s1.c_str());
    cout << num << endl;
    
    string s2 = "3.14";
    double dbl = atof(s2.c_str());
    cout << dbl << endl;
    
    // Number to string (using stringstream)
    stringstream ss;
    ss << 42 << " " << 3.14 << " " << true;
    string result = ss.str();
    cout << result << endl;  // 42 3.14 1
    
    // Reverse conversion
    stringstream ss2("100 200 300");
    int a, b, c;
    ss2 >> a >> b >> c;
    cout << a << " " << b << " " << c << endl;  // 100 200 300
    
    return 0;
}
\end{lstlisting}

\hypertarget{string-tokenization}{%
\section{5.3 String Tokenization}\label{string-tokenization}}

\begin{lstlisting}[language={C++}]
#include <iostream>
#include <string>
#include <cstring>
using namespace std;

int main() {
    string line = "apple,banana,orange,grape";
    
    // Using stringstream and getline
    stringstream ss(line);
    string token;
    
    while (getline(ss, token, ',')) {
        cout << token << endl;
    }
    // Output: apple, banana, orange, grape
    
    // Using strtok (C style)
    char str[] = "hello world how are you";
    char* ptr = strtok(str, " ");
    
    while (ptr != NULL) {
        cout << ptr << endl;
        ptr = strtok(NULL, " ");
    }
    // Output: hello, world, how, are, you
    
    return 0;
}
\end{lstlisting}

\begin{center}\rule{0.5\linewidth}{0.5pt}\end{center}

\hypertarget{file-io-advanced}{%
\chapter{FILE I/O ADVANCED}\label{file-io-advanced}}

\hypertarget{binary-file-io}{%
\section{14.1 Binary File I/O}\label{binary-file-io}}

\begin{lstlisting}[language={C++}]
#include <iostream>
#include <fstream>
using namespace std;

int main() {
    // Write binary data
    ofstream outfile("data.bin", ios::binary);
    
    int numbers[] = {10, 20, 30, 40, 50};
    outfile.write((char*)numbers, sizeof(numbers));
    outfile.close();
    
    // Read binary data
    ifstream infile("data.bin", ios::binary);
    
    int buffer[5];
    infile.read((char*)buffer, sizeof(buffer));
    
    for (int i = 0; i < 5; i++) {
        cout << buffer[i] << " ";
    }
    cout << endl;
    
    infile.close();
    
    return 0;
}
\end{lstlisting}

\hypertarget{stream-positioning}{%
\section{14.2 Stream Positioning}\label{stream-positioning}}

\begin{lstlisting}[language={C++}]
#include <iostream>
#include <fstream>
using namespace std;

int main() {
    // Write to file
    ofstream outfile("test.txt");
    outfile << "0123456789";
    outfile.close();
    
    // Read with positioning
    ifstream infile("test.txt");
    
    // Tell position
    cout << "Current position: " << infile.tellg() << endl;
    
    // Seek to position
    infile.seekg(5);
    char c;
    infile.get(c);
    cout << "Character at position 5: " << c << endl;  // '5'
    
    // Seek from end
    infile.seekg(-3, ios::end);
    infile.get(c);
    cout << "Third from end: " << c << endl;  // '7'
    
    infile.close();
    
    return 0;
}
\end{lstlisting}

\begin{center}\rule{0.5\linewidth}{0.5pt}\end{center}

\hypertarget{professional-notes-chapter-47-stdstring}{%
\chapter{Professional Notes: Chapter 47:
std::string}\label{professional-notes-chapter-47-stdstring}}

Section 47.1: Tokenize Section 47.2: Conversion to (const) char* Section
47.3: Using the std::string\_view class Section 47.4: Conversion to
std::wstring Section 47.5: Lexicographical comparison Section 47.6:
Trimming characters at start/end Section 47.7: String replacement
Section 47.8: Converting to std::string Section 47.9: Splitting Section
47.10: Accessing a character Section 47.11: Checking if a string is a
prex of another Section 47.12: Looping through each character Section
47.13: Conversion to integers/oating point types Section 47.14:
Concatenation Section 47.15: Converting between character encodings
Section 47.16: Finding character(s) in a string

\hypertarget{professional-notes-chapter-49-stdvector}{%
\chapter{Professional Notes: Chapter 49:
std::vector}\label{professional-notes-chapter-49-stdvector}}

Section 49.1: Accessing Elements Section 49.2: Initializing a
std::vector Section 49.3: Deleting Elements Section 49.4: Iterating Over
std::vector Section 49.5: vector: The Exception To So Many, So Many
Rules Section 49.6: Inserting Elements Section 49.7: Using std::vector
as a C array Section 49.8: Finding an Element in std::vector Section
49.9: Concatenating Vectors Section 49.10: Matrices Using Vectors
Section 49.11: Using a Sorted Vector for Fast Element Lookup Section
49.12: Reducing the Capacity of a Vector Section 49.13: Vector size and
capacity Section 49.14: Iterator/Pointer Invalidation Section 49.15:
Find max and min Element and Respective Index in a Vector Section 49.16:
Converting an array to std::vector Section 49.17: Functions Returning
Large Vectors

\hypertarget{professional-notes-chapter-50-stdmap}{%
\chapter{Professional Notes: Chapter 50:
std::map}\label{professional-notes-chapter-50-stdmap}}

Section 50.1: Accessing elements Section 50.2: Inserting elements
Section 50.3: Searching in std::map or in std::multimap Section 50.4:
Initializing a std::map or std::multimap Section 50.5: Checking number
of elements Section 50.6: Types of Maps Section 50.7: Deleting elements
Section 50.8: Iterating over std::map or std::multimap Section 50.9:
Creating std::map with user-dened types as key

\hypertarget{professional-notes-chapter-62-standard-library-algorithms}{%
\chapter{Professional Notes: Chapter 62: Standard Library
Algorithms}\label{professional-notes-chapter-62-standard-library-algorithms}}

Section 62.1: std::next\_permutation Section 62.2: std::for\_each
Section 62.3: std::accumulate Section 62.4: std::nd Section 62.5:
std::min\_element Section 62.6: std::nd\_if Section 62.7: Using
std::nth\_element To Find The Median (Or Other Quantiles) Section 62.8:
std::count Section 62.9: std::count\_if

\hypertarget{professional-notes-chapter-67-sorting}{%
\chapter{Professional Notes: Chapter 67:
Sorting}\label{professional-notes-chapter-67-sorting}}

Section 67.1: Sorting and sequence containers Section 67.2: sorting with
std::map (ascending and descending) Section 67.3: Sorting sequence
containers by overloaded less operator Section 67.4: Sorting sequence
containers using compare function Section 67.5: Sorting sequence
containers using lambda expressions (C++11) Section 67.6: Sorting
built-in arrays Section 67.7: Sorting sequence containers with specifed
ordering

\hypertarget{professional-notes-chapter-43-c-containers}{%
\chapter{Professional Notes: Chapter 43: C++
Containers}\label{professional-notes-chapter-43-c-containers}}

C++ containers store a collection of elements. Containers include
vectors, lists, maps, etc. Using Templates, C++ containers contain
collections of primitives (e.g.~ints) or custom classes (e.g.~MyClass).
Section 43.1: C++ Containers Flowchart Choosing which C++ Container to
use can be tricky, so here's a simple owchart to help decide which
Container is right for the job. This owchart was based on Mikael
Persson's post. This little graphic in the owchart is from Megan Hopkins

\hypertarget{professional-notes-chapter-47-stdstring-1}{%
\chapter{Professional Notes: Chapter 47:
std::string}\label{professional-notes-chapter-47-stdstring-1}}

Strings are objects that represent sequences of characters. The standard
string class provides a simple, safe and versatile alternative to using
explicit arrays of chars when dealing with text and other sequences of
characters. The C++ string class is part of the std namespace and was
standardized in 1998. Section 47.1: Tokenize Listed from least expensive
to most expensive at run-time: 1. std::strtok is the cheapest standard
provided tokenization method, it also allows the delimiter to be modied
between tokens, but it incurs 3 diculties with modern C++: std::strtok
cannot be used on multiple strings at the same time (though some
implementations do extend to support this, such as: strtok\_s) For the
same reason std::strtok cannot be used on multiple threads
simultaneously (this may however be implementation dened, for example:
Visual Studio's implementation is thread safe) Calling std::strtok
modies the std::string it is operating on, so it cannot be used on const
strings, const char\emph{s, or literal strings, to tokenize any of these
with std::strtok or to operate on a std::string who's contents need to
be preserved, the input would have to be copied, then the copy could be
operated on Generally any of these options cost will be hidden in the
allocation cost of the tokens, but if the cheapest algorithm is required
and std::strtok's diculties are not overcomable consider a hand-spun
solution. // String to tokenize std::string str\{ ``The quick brown
fox'' \}; // Vector to store tokens vector tokens; for (auto i =
strtok(\&str{[}0{]}, " ``); i != NULL; i = strtok(NULL,'' ``))
tokens.push\_back(i); Live Example 2. The std::istream\_iterator uses
the stream's extraction operator iteratively. If the input std::string
is white-space delimited this is able to expand on the std::strtok
option by eliminating its diculties, allowing inline tokenization
thereby supporting the generation of a const vector, and by adding
support for multiple delimiting white-space character: // String to
tokenize const std::string str(''The quick \tbrown \nfox``);
std::istringstream is(str); // Vector to store tokens const std::vector
tokens = std::vector( std::istream\_iterator(is),
std::istream\_iterator()); Live Example 3. The
std::regex\_token\_iterator uses a std::regex to iteratively tokenize.
It provides for a more exible delimiter denition. For example,
non-delimited commas and white-space: Version C++11 // String to
tokenize const std::string str\{''The ,qu\textbackslash,ick ,\tbrown,
fox" \}; const std::regex re\{
``\textbackslash s\emph{((?:{[}\^{}\textbackslash\textbackslash\textbackslash\textbackslash,{]}\textbar\textbackslash\textbackslash.)}?)\textbackslash s\emph{(?:,\textbar\$)"
\}; // Vector to store tokens const std::vector tokens\{
std::sregex\_token\_iterator(str.begin(), str.end(), re, 1),
std::sregex\_token\_iterator() \}; Live Example See the
regex\_token\_iterator Example for more details. Section 47.2:
Conversion to (const) char} In order to get const char* access to the
data of a std::string you can use the string's c\_str() member function.
Keep in mind that the pointer is only valid as long as the std::string
object is within scope and remains unchanged, that means that only const
methods may be called on the object. Version C++17 The data() member
function can be used to obtain a modiable char\emph{, which can be used
to manipulate the std::string object's data. Version C++11 A modiable
char} can also be obtained by taking the address of the rst character:
\&s{[}0{]}. Within C++11, this is guaranteed to yield a well-formed,
null-terminated string. Note that \&s{[}0{]} is well-formed even if s is
empty, whereas \&s.front() is undened if s is empty. Version C++11
std::string str(''This is a string.``); const char* cstr = str.c\_str();
// cstr points to:''This is a string.\textbackslash0" const char} data =
str.data(); // data points to: ``This is a string.\textbackslash0''
std::string str(``This is a string.''); // Copy the contents of str to
untie lifetime from the std::string object
std::unique\_ptr\textless char {[}{]}\textgreater{} cstr =
std::make\_unique\textless char{[}{]}\textgreater(str.size() + 1); //
Alternative to the line above (no exception safety): // char*
cstr\_unsafe = new char{[}str.size() + 1{]}; std::copy(str.data(),
str.data() + str.size(), cstr); cstr{[}str.size(){]} =
`\textbackslash0'; // A null-terminator needs to be added //
delete{[}{]} cstr\_unsafe; std::cout \textless\textless{} cstr.get();
Section 47.3: Using the std::string\_view class Version C++17 C++17
introduces std::string\_view, which is simply a non-owning range of
const chars, implementable as either a pair of pointers or a pointer and
a length. It is a superior parameter type for functions that requires
non- modiable string data. Before C++17, there were three options for
this: void foo(std::string const\& s); // pre-C++17, single argument,
could incur  // allocation if caller's data was not in a string //
(e.g.~string literal or vector ) void foo(const char* s, size\_t len);
// pre-C++17, two arguments, have to pass them // both everywhere void
foo(const char* s); // pre-C++17, single argument, but need to call //
strlen() template void foo(StringT const\& s); // pre-C++17, caller can
pass arbitrary char data // provider, but now foo() has to live in a
header All of these can be replaced with: void foo(std::string\_view s);
// post-C++17, single argument, tighter coupling // zero copies
regardless of how caller is storing // the data Note that
std::string\_view cannot modify its underlying data. string\_view is
useful when you want to avoid unnecessary copies. It oers a useful
subset of the functionality that std::string does, although some of the
functions behave dierently: std::string str = ``lllloooonnnngggg
sssstttrrriiinnnggg''; //A really long string //Bad way -
`string::substr' returns a new string (expensive if the string is long)
std::cout \textless\textless{} str.substr(15, 10) \textless\textless{}
`\n'; //Good way - No copies are created! std::string\_view view = str;
// string\_view::substr returns a new string\_view std::cout
\textless\textless{} view.substr(15, 10) \textless\textless{} `\n';
Section 47.4: Conversion to std::wstring In C++, sequences of characters
are represented by specializing the std::basic\_string class with a
native character type. The two major collections dened by the standard
library are std::string and std::wstring: std::string is built with
elements of type char std::wstring is built with elements of type
wchar\_t To convert between the two types, use wstring\_convert:
\#include \#include \#include std::string input\_str = ``this is a
-string-, which is a sequence based on the -char- type.''; std::wstring
input\_wstr = L``this is a -wide- string, which is based on the
-wchar\_t- type.''; // conversion std::wstring str\_turned\_to\_wstr =
std::wstring\_convert\textless std::codecvt\_utf8\textgreater().from\_bytes(input\_str);
std::string wstr\_turned\_to\_str =
std::wstring\_convert\textless std::codecvt\_utf8\textgreater().to\_bytes(input\_wstr);
In order to improve usability and/or readability, you can dene functions
to perform the conversion: \#include \#include \#include using
convert\_t = std::codecvt\_utf8;
std::wstring\_convert\textless convert\_t, wchar\_t\textgreater{}
strconverter; std::string to\_string(std::wstring wstr) \{ return
strconverter.to\_bytes(wstr); \} std::wstring to\_wstring(std::string
str) \{ return strconverter.from\_bytes(str); \} Sample usage:
std::wstring a\_wide\_string = to\_wstring(``Hello World!''); That's
certainly more readable than
std::wstring\_convert\textless std::codecvt\_utf8\textgreater().from\_bytes(``Hello
World!''). Please note that char and wchar\_t do not imply encoding, and
gives no indication of size in bytes. For instance, wchar\_t is commonly
implemented as a 2-bytes data type and typically contains UTF-16 encoded
data under Windows (or UCS-2 in versions prior to Windows 2000) and as a
4-bytes data type encoded using UTF-32 under Linux. This is in contrast
with the newer types char16\_t and char32\_t, which were introduced in
C++11 and are guaranteed to be large enough to hold any UTF16 or UTF32
``character'' (or more precisely, code point) respectively. Section
47.5: Lexicographical comparison Two std::strings can be compared
lexicographically using the operators ==, !=, \textless, \textless=,
\textgreater, and \textgreater=: std::string str1 = ``Foo''; std::string
str2 = ``Bar''; assert(!(str1 \textless{} str2)); assert(str
\textgreater{} str2); assert(!(str1 \textless= str2)); assert(str1
\textgreater= str2); assert(!(str1 == str2)); assert(str1 != str2); All
these functions use the underlying std::string::compare() method to
perform the comparison, and return for convenience boolean values. The
operation of these functions may be interpreted as follows, regardless
of the actual implementation: operator==: If str1.length() ==
str2.length() and each character pair matches, then returns true,
otherwise returns false. operator!=: If str1.length() != str2.length()
or one character pair doesn't match, returns true, otherwise it returns
false. operator\textless{} or operator\textgreater: Finds the rst
dierent character pair, compares them then returns the boolean result.
operator\textless= or operator\textgreater=: Finds the rst dierent
character pair, compares them then returns the boolean result. Note: The
term character pair means the corresponding characters in both strings
of the same positions. For better understanding, if two example strings
are str1 and str2, and their lengths are n and m respectively, then
character pairs of both strings means each str1{[}i{]} and str2{[}i{]}
pairs where i = 0, 1, 2, \ldots, max(n,m). If for any i where the
corresponding character does not exist, that is, when i is greater than
or equal to n or m, it would be considered as the lowest value. Here is
an example of using \textless: std::string str1 = ``Barr''; std::string
str2 = ``Bar''; assert(str2 \textless{} str1); The steps are as follows:
1. 2. 3. 4. Compare the rst characters, `B' == `B' - move on. Compare
the second characters, `a' == `a' - move on. Compare the third
characters, `r' == `r' - move on. The str2 range is now exhausted, while
the str1 range still has characters. Thus, str2 \textless{} str1.
Section 47.6: Trimming characters at start/end This example requires the
headers , , and . Version C++11 To trim a sequence or string means to
remove all leading and trailing elements (or characters) matching a
certain predicate. We rst trim the trailing elements, because it doesn't
involve moving any elements, and then trim the leading elements. Note
that the generalizations below work for all types of std::basic\_string
(e.g.~std::string and std::wstring), and accidentally also for sequence
containers (e.g.~std::vector and std::list). template \textless typename
Sequence, // any basic\_string, vector, list etc. typename
Pred\textgreater{} // a predicate on the element (character) type
Sequence\& trim(Sequence\& seq, Pred pred) \{ return
trim\_start(trim\_end(seq, pred), pred); \} Trimming the trailing
elements involves nding the last element not matching the predicate, and
erasing from there on: template \textless typename Sequence, typename
Pred\textgreater{} Sequence\& trim\_end(Sequence\& seq, Pred pred) \{
auto last = std::find\_if\_not(seq.rbegin(), seq.rend(), pred);
seq.erase(last.base(), seq.end()); return seq; \} Trimming the leading
elements involves nding the rst element not matching the predicate and
erasing up to there: template \textless typename Sequence, typename
Pred\textgreater{} Sequence\& trim\_start(Sequence\& seq, Pred pred) \{
auto first = std::find\_if\_not(seq.begin(), seq.end(), pred);
seq.erase(seq.begin(), first); return seq; \} To specialize the above
for trimming whitespace in a std::string we can use the std::isspace()
function as a predicate: std::string\& trim(std::string\& str, const
std::locale\& loc = std::locale()) \{ return trim(str,
\href{const\%20char\%20c}{\&loc}\{ return std::isspace(c, loc); \}); \}
std::string\& trim\_start(std::string\& str, const std::locale\& loc =
std::locale()) \{ return trim\_start(str,
\href{const\%20char\%20c}{\&loc}\{ return std::isspace(c, loc); \}); \}
std::string\& trim\_end(std::string\& str, const std::locale\& loc =
std::locale()) \{ return trim\_end(str,
\href{const\%20char\%20c}{\&loc}\{ return std::isspace(c, loc); \}); \}
Similarly, we can use the std::iswspace() function for std::wstring etc.
If you wish to create a new sequence that is a trimmed copy, then you
can use a separate function: template \textless typename Sequence,
typename Pred\textgreater{} Sequence trim\_copy(Sequence seq, Pred pred)
\{ // NOTE: passing seq by value trim(seq, pred); return seq; \} Section
47.7: String replacement Replace by position To replace a portion of a
std::string you can use the method replace from std::string. replace has
a lot of useful overloads: //Define string std::string str = ``Hello
foo, bar and world!''; std::string alternate = ``Hello foobar''; //1)
str.replace(6, 3, ``bar''); //``Hello bar, bar and world!'' //2)
str.replace(str.begin() + 6, str.end(), ``nobody!''); //``Hello
nobody!'' //3) str.replace(19, 5, alternate, 6, 6); //``Hello foo, bar
and foobar!'' Version C++14 //4) str.replace(19, 5, alternate, 6);
//``Hello foo, bar and foobar!'' //5) str.replace(str.begin(),
str.begin() + 5, str.begin() + 6, str.begin() + 9); //``foo foo, bar and
world!'' //6) str.replace(0, 5, 3, `z'); //``zzz foo, bar and world!''
//7) str.replace(str.begin() + 6, str.begin() + 9, 3, `x'); //``Hello
xxx, bar and world!'' Version C++11 //8) str.replace(str.begin(),
str.begin() + 5, \{ `x', `y', `z' \}); //``xyz foo, bar and world!''
Replace occurrences of a string with another string Replace only the rst
occurrence of replace with with in str: std::string
replaceString(std::string str, const std::string\& replace, const
std::string\& with)\{ std::size\_t pos = str.find(replace); if (pos !=
std::string::npos) str.replace(pos, replace.length(), with); return str;
\} Replace all occurrence of replace with with in str: std::string
replaceStringAll(std::string str, const std::string\& replace, const
std::string\& with) \{ if(!replace.empty()) \{ std::size\_t pos = 0;
while ((pos = str.find(replace, pos)) != std::string::npos) \{
str.replace(pos, replace.length(), with); pos += with.length(); \} \}
return str; \} Section 47.8: Converting to std::string
std::ostringstream can be used to convert any streamable type to a
string representation, by inserting the object into a
std::ostringstream object (with the stream insertion operator
\textless\textless) and then converting the whole std::ostringstream to
a std::string. For int for instance: \#include int main() \{ int val =
4; std::ostringstream str; str \textless\textless{} val; std::string
converted = str.str(); return 0; \} Writing your own conversion
function, the simple: template std::string toString(const T\& x) \{
std::ostringstream ss; ss \textless\textless{} x; return ss.str(); \}
works but isn't suitable for performance critical code. User-dened
classes may implement the stream insertion operator if desired:
std::ostream operator\textless\textless( std::ostream\& out, const A\& a
) \{ // write a string representation of a to out return out; \} Version
C++11 Aside from streams, since C++11 you can also use the
std::to\_string (and std::to\_wstring) function which is overloaded for
all fundamental types and returns the string representation of its
parameter. std::string s = to\_string(0x12f3); // after this the string
s contains ``4851'' Section 47.9: Splitting Use std::string::substr to
split a string. There are two variants of this member function. The rst
takes a starting position from which the returned substring should
begin. The starting position must be valid in the range (0,
str.length(){]}: std::string str = ``Hello foo, bar and world!'';
std::string newstr = str.substr(11); // ``bar and world!'' The second
takes a starting position and a total length of the new substring.
Regardless of the length, the substring will never go past the end of
the source string: std::string str = ``Hello foo, bar and world!'';
std::string newstr = str.substr(15, 3); // ``and'' Note that you can
also call substr with no arguments, in this case an exact copy of the
string is returned std::string str = ``Hello foo, bar and world!'';
std::string newstr = str.substr(); // ``Hello foo, bar and world!''
Section 47.10: Accessing a character There are several ways to extract
characters from a std::string and each is subtly dierent. std::string
str(``Hello world!''); operator\href{n}{} Returns a reference to the
character at index n. std::string::operator{[}{]} is not bounds-checked
and does not throw an exception. The caller is responsible for asserting
that the index is within the range of the string: char c = str{[}6{]};
// `w' at(n) Returns a reference to the character at index n.
std::string::at is bounds checked, and will throw std::out\_of\_range if
the index is not within the range of the string: char c = str.at(7); //
`o' Version C++11 Note: Both of these examples will result in undened
behavior if the string is empty. front() Returns a reference to the rst
character: char c = str.front(); // `H' back() Returns a reference to
the last character: char c = str.back(); // `!' Section 47.11: Checking
if a string is a prex of another Version C++14 In C++14, this is easily
done by std::mismatch which returns the rst mismatching pair from two
ranges: std::string prefix = ``foo''; std::string string = ``foobar'';
bool isPrefix = std::mismatch(prefix.begin(), prefix.end(),
string.begin(), string.end()).first == prefix.end(); Note that a
range-and-a-half version of mismatch() existed prior to C++14, but this
is unsafe in the case that the second string is the shorter of the two.
Version \textless{} C++14 We can still use the range-and-a-half version
of std::mismatch(), but we need to rst check that the rst string is at
most as big as the second: bool isPrefix = prefix.size() \textless=
string.size() \&\& std::mismatch(prefix.begin(), prefix.end(),
string.begin(), string.end()).first == prefix.end(); Version C++17 With
std::string\_view, we can write the direct comparison we want without
having to worry about allocation overhead or making copies: bool
isPrefix(std::string\_view prefix, std::string\_view full) \{ return
prefix == full.substr(0, prefix.size()); \} Section 47.12: Looping
through each character Version C++11 std::string supports iterators, and
so you can use a ranged based loop to iterate through each character:
std::string str = ``Hello World!''; for (auto c : str) std::cout
\textless\textless{} c; You can use a ``traditional'' for loop to loop
through every character: std::string str = ``Hello World!''; for
(std::size\_t i = 0; i \textless{} str.length(); ++i) std::cout
\textless\textless{} str{[}i{]}; Section 47.13: Conversion to
integers/oating point types A std::string containing a number can be
converted into an integer type, or a oating point type, using conversion
functions. Note that all of these functions stop parsing the input
string as soon as they encounter a non-numeric character, so ``123abc''
will be converted into 123. The std::ato* family of functions converts
C-style strings (character arrays) to integer or oating-point types:
std::string ten = ``10''; double num1 = std::atof(ten.c\_str()); int
num2 = std::atoi(ten.c\_str()); long num3 = std::atol(ten.c\_str());
Version C++11 long long num4 = std::atoll(ten.c\_str()); However, use of
these functions is discouraged because they return 0 if they fail to
parse the string. This is bad because 0 could also be a valid result, if
for example the input string was ``0'', so it is impossible to determine
if the conversion actually failed. The newer std::sto* family of
functions convert std::strings to integer or oating-point types, and
throw exceptions if they could not parse their input. You should use
these functions if possible: Version C++11 std::string ten = ``10''; int
num1 = std::stoi(ten); long num2 = std::stol(ten); long long num3 =
std::stoll(ten); float num4 = std::stof(ten); double num5 =
std::stod(ten); long double num6 = std::stold(ten); Furthermore, these
functions also handle octal and hex strings unlike the std::ato* family.
The second parameter is a pointer to the rst unconverted character in
the input string (not illustrated here), and the third parameter is the
base to use. 0 is automatic detection of octal (starting with 0) and hex
(starting with 0x or 0X), and any other value is the base to use
std::string ten = ``10''; std::string ten\_octal = ``12''; std::string
ten\_hex = ``0xA''; int num1 = std::stoi(ten, 0, 2); // Returns 2 int
num2 = std::stoi(ten\_octal, 0, 8); // Returns 10 long num3 =
std::stol(ten\_hex, 0, 16); // Returns 10 long num4 =
std::stol(ten\_hex); // Returns 0 long num5 = std::stol(ten\_hex, 0, 0);
// Returns 10 as it detects the leading 0x Section 47.14: Concatenation
You can concatenate std::strings using the overloaded + and +=
operators. Using the + operator: std::string hello = ``Hello'';
std::string world = ``world''; std::string helloworld = hello + world;
// ``Helloworld'' Using the += operator: std::string hello = ``Hello'';
std::string world = ``world''; hello += world; // ``Helloworld'' You can
also append C strings, including string literals: std::string hello =
``Hello''; std::string world = ``world''; const char *comma = ``,'';
std::string newhelloworld = hello + comma + world + ``!''; // ``Hello,
world!'' You can also use push\_back() to push back individual chars:
std::string s = ``a, b,''; s.push\_back(`c'); // ``a, b, c'' There is
also append(), which is pretty much like +=: std::string app = ``test
and''; app.append(``test''); // ``test and test'' Section 47.15:
Converting between character encodings Converting between encodings is
easy with C++11 and most compilers are able to deal with it in a
cross-platform manner through and headers. \#include \#include \#include
\#include using namespace std; int main() \{ // converts between wstring
and utf8 string
wstring\_convert\textless codecvt\_utf8\_utf16\textgreater{}
wchar\_to\_utf8; // converts between u16string and utf8 string
wstring\_convert\textless codecvt\_utf8\_utf16, char16\_t\textgreater{}
utf16\_to\_utf8; wstring wstr = L``foobar''; string utf8str =
wchar\_to\_utf8.to\_bytes(wstr); wstring wstr2 =
wchar\_to\_utf8.from\_bytes(utf8str); wcout \textless\textless{} wstr
\textless\textless{} endl; cout \textless\textless{} utf8str
\textless\textless{} endl; wcout \textless\textless{} wstr2
\textless\textless{} endl; u16string u16str = u``foobar''; string
utf8str2 = utf16\_to\_utf8.to\_bytes(u16str); u16string u16str2 =
utf16\_to\_utf8.from\_bytes(utf8str2); return 0; \} Mind that Visual
Studio 2015 provides supports for these conversion but a bug in their
library implementation requires to use a dierent template for
wstring\_convert when dealing with char16\_t: using utf16\_char =
unsigned short; wstring\_convert\textless codecvt\_utf8\_utf16,
utf16\_char\textgreater{} conv\_utf8\_utf16; void
strings::utf16\_to\_utf8(const std::u16string\& utf16, std::string\&
utf8) \{ std::basic\_string tmp; tmp.resize(utf16.length());
std::copy(utf16.begin(), utf16.end(), tmp.begin()); utf8 =
conv\_utf8\_utf16.to\_bytes(tmp); \} void
strings::utf8\_to\_utf16(const std::string\& utf8, std::u16string\&
utf16) \{ std::basic\_string tmp = conv\_utf8\_utf16.from\_bytes(utf8);
utf16.clear(); utf16.resize(tmp.length()); std::copy(tmp.begin(),
tmp.end(), utf16.begin()); \} Section 47.16: Finding character(s) in a
string To nd a character or another string, you can use
std::string::find. It returns the position of the rst character of the
rst match. If no matches were found, the function returns
std::string::npos std::string str = ``Curiosity killed the cat''; auto
it = str.find(``cat''); if (it != std::string::npos) std::cout
\textless\textless{} ``Found at position:'' \textless\textless{} it
\textless\textless{} `\n'; else std::cout \textless\textless{} ``Not
found!\n''; Found at position: 21 The search opportunities are further
expanded by the following functions: find\_first\_of // Find first
occurrence of characters find\_first\_not\_of // Find first absence of
characters find\_last\_of // Find last occurrence of characters
find\_last\_not\_of // Find last absence of characters These functions
can allow you to search for characters from the end of the string, as
well as nd the negative case (ie. characters that are not in the
string). Here is an example: std::string str = ``dog dog cat cat'';
std::cout \textless\textless{} ``Found at position:''
\textless\textless{} str.find\_last\_of(``gzx'') \textless\textless{}
`\n'; Found at position: 6 Note: Be aware that the above functions do
not search for substrings, but rather for characters contained in the
search string. In this case, the last occurrence of `g' was found at
position 6 (the other characters weren't found).

\hypertarget{professional-notes-chapter-49-stdvector-1}{%
\chapter{Professional Notes: Chapter 49:
std::vector}\label{professional-notes-chapter-49-stdvector-1}}

A vector is a dynamic array with automatically handled storage. The
elements in a vector can be accessed just as eciently as those in an
array with the advantage being that vectors can dynamically change in
size. In terms of storage the vector data is (usually) placed in
dynamically allocated memory thus requiring some minor overhead;
conversely C-arrays and std::array use automatic storage relative to the
declared location and thus do not have any overhead. Section 49.1:
Accessing Elements There are two primary ways of accessing elements in a
std::vector index-based access iterators Index-based access: This can be
done either with the subscript operator {[}{]}, or the member function
at(). Both return a reference to the element at the respective position
in the std::vector (unless it's a vector), so that it can be read as
well as modied (if the vector is not const). {[}{]} and at() dier in
that {[}{]} is not guaranteed to perform any bounds checking, while at()
does. Accessing elements where index \textless{} 0 or index
\textgreater= size is undened behavior for {[}{]}, while at() throws a
std::out\_of\_range exception. Note: The examples below use C++11-style
initialization for clarity, but the operators can be used with all
versions (unless marked C++11). Version C++11 std::vector v\{ 1, 2, 3
\}; // using {[}{]} int a = v{[}1{]}; // a is 2 v{[}1{]} = 4; // v now
contains \{ 1, 4, 3 \} // using at() int b = v.at(2); // b is 3 v.at(2)
= 5; // v now contains \{ 1, 4, 5 \} int c = v.at(3); // throws
std::out\_of\_range exception Because the at() method performs bounds
checking and can throw exceptions, it is slower than {[}{]}. This makes
{[}{]} preferred code where the semantics of the operation guarantee
that the index is in bounds. In any case, accesses to elements of
vectors are done in constant time. That means accessing to the rst
element of the vector has the same cost (in time) of accessing the
second element, the third element and so on. For example, consider this
loop for (std::size\_t i = 0; i \textless{} v.size(); ++i) \{ v{[}i{]} =
1; \} Here we know that the index variable i is always in bounds, so it
would be a waste of CPU cycles to check that i is in bounds for every
call to operator{[}{]}. The front() and back() member functions allow
easy reference access to the rst and last element of the vector,
respectively. These positions are frequently used, and the special
accessors can be more readable than their alternatives using {[}{]}:
std::vector v\{ 4, 5, 6 \}; // In pre-C++11 this is more verbose int a =
v.front(); // a is 4, v.front() is equivalent to v{[}0{]} v.front() = 3;
// v now contains \{3, 5, 6\} int b = v.back(); // b is 6, v.back() is
equivalent to v{[}v.size() - 1{]} v.back() = 7; // v now contains \{3,
5, 7\} Note: It is undened behavior to invoke front() or back() on an
empty vector. You need to check that the container is not empty using
the empty() member function (which checks if the container is empty)
before calling front() or back(). A simple example of the use of
`empty()' to test for an empty vector follows: int main () \{
std::vector v; int sum (0); for (int i=1;i\textless=10;i++)
v.push\_back(i);//create and initialize the vector while
(!v.empty())//loop through until the vector tests to be empty \{ sum +=
v.back();//keep a running total v.pop\_back();//pop out the element
which removes it from the vector \} std::cout \textless\textless{}
``total:'' \textless\textless{} sum \textless\textless{} `\n';//output
the total to the user return 0; \} The example above creates a vector
with a sequence of numbers from 1 to 10. Then it pops the elements of
the vector out until the vector is empty (using `empty()') to prevent
undened behavior. Then the sum of the numbers in the vector is
calculated and displayed to the user. Version C++11 The data() method
returns a pointer to the raw memory used by the std::vector to
internally store its elements. This is most often used when passing the
vector data to legacy code that expects a C-style array. std::vector v\{
1, 2, 3, 4 \}; // v contains \{1, 2, 3, 4\} int* p = v.data(); // p
points to 1 \emph{p = 4; // v now contains \{4, 2, 3, 4\} ++p; // p
points to 2 }p = 3; // v now contains \{4, 3, 3, 4\} p{[}1{]} = 2; // v
now contains \{4, 3, 2, 4\} \emph{(p + 2) = 1; // v now contains \{4, 3,
2, 1\} Version \textless{} C++11 Before C++11, the data() method can be
simulated by calling front() and taking the address of the returned
value: std::vector v(4); int} ptr = \&(v.front()); // or \&v{[}0{]} This
works because vectors are always guaranteed to store their elements in
contiguous memory locations, assuming the contents of the vector
doesn't override unary operator\&. If it does, you'll have to
re-implement std::addressof in pre-C++11. It also assumes that the
vector isn't empty. Iterators: Iterators are explained in more detail in
the example ``Iterating over std::vector'' and the article Iterators. In
short, they act similarly to pointers to the elements of the vector:
Version C++11 std::vector v\{ 4, 5, 6 \}; auto it = v.begin(); int i =
\emph{it; // i is 4 ++it; i = }it; // i is 5 \emph{it = 6; // v contains
\{ 4, 6, 6 \} auto e = v.end(); // e points to the element after the end
of v. It can be // used to check whether an iterator reached the end of
the vector: ++it; it == v.end(); // false, it points to the element at
position 2 (with value 6) ++it; it == v.end(); // true It is consistent
with the standard that a std::vector's iterators actually be T}s, but
most standard libraries do not do this. Not doing this both improves
error messages, catches non-portable code, and can be used to instrument
the iterators with debugging checks in non-release builds. Then, in
release builds, the class wrapping around the underlying pointer is
optimized away. You can persist a reference or a pointer to an element
of a vector for indirect access. These references or pointers to
elements in the vector remain stable and access remains dened unless you
add/remove elements at or before the element in the vector, or you cause
the vector capacity to change. This is the same as the rule for
invalidating iterators. Version C++11 std::vector v\{ 1, 2, 3 \}; int* p
= v.data() + 1; // p points to 2 v.insert(v.begin(), 0); // p is now
invalid, accessing \emph{p is a undefined behavior. p = v.data() + 1; //
p points to 1 v.reserve(10); // p is now invalid, accessing }p is a
undefined behavior. p = v.data() + 1; // p points to 1
v.erase(v.begin()); // p is now invalid, accessing *p is a undefined
behavior. Section 49.2: Initializing a std::vector A std::vector can be
initialized in several ways while declaring it: Version C++11
std::vector v\{ 1, 2, 3 \}; // v becomes \{1, 2, 3\} // Different from
std::vector v(3, 6) std::vector v\{ 3, 6 \}; // v becomes \{3, 6\} //
Different from std::vector v\{3, 6\} in C++11 std::vector v(3, 6); // v
becomes \{6, 6, 6\} std::vector v(4); // v becomes \{0, 0, 0, 0\} A
vector can be initialized from another container in several ways: Copy
construction (from another vector only), which copies data from v2:
std::vector v(v2); std::vector v = v2; Version C++11 Move construction
(from another vector only), which moves data from v2: std::vector
v(std::move(v2)); std::vector v = std::move(v2); Iterator (range)
copy-construction, which copies elements into v: // from another vector
std::vector v(v2.begin(), v2.begin() + 3); // v becomes \{v2{[}0{]},
v2{[}1{]}, v2{[}2{]}\} // from an array int z{[}{]} = \{ 1, 2, 3, 4 \};
std::vector v(z, z + 3); // v becomes \{1, 2, 3\} // from a list
std::list list1\{ 1, 2, 3 \}; std::vector v(list1.begin(), list1.end());
// v becomes \{1, 2, 3\} Version C++11 Iterator move-construction, using
std::make\_move\_iterator, which moves elements into v: // from another
vector std::vector v(std::make\_move\_iterator(v2.begin()),
std::make\_move\_iterator(v2.end()); // from a list std::list list1\{ 1,
2, 3 \}; std::vector v(std::make\_move\_iterator(list1.begin()),
std::make\_move\_iterator(list1.end())); With the help of the assign()
member function, a std::vector can be reinitialized after its
construction: v.assign(4, 100); // v becomes \{100, 100, 100, 100\}
v.assign(v2.begin(), v2.begin() + 3); // v becomes \{v2{[}0{]},
v2{[}1{]}, v2{[}2{]}\} int z{[}{]} = \{ 1, 2, 3, 4 \}; v.assign(z + 1, z
+ 4); // v becomes \{2, 3, 4\} Section 49.3: Deleting Elements Deleting
the last element: std::vector v\{ 1, 2, 3 \}; v.pop\_back(); // v
becomes \{1, 2\} Deleting all elements: std::vector v\{ 1, 2, 3 \};
v.clear(); // v becomes an empty vector Deleting element by index:
std::vector v\{ 1, 2, 3, 4, 5, 6 \}; v.erase(v.begin() + 3); // v
becomes \{1, 2, 3, 5, 6\} Note: For a vector deleting an element which
is not the last element, all elements beyond the deleted element have to
be copied or moved to ll the gap, see the note below and std::list.
Deleting all elements in a range: std::vector v\{ 1, 2, 3, 4, 5, 6 \};
v.erase(v.begin() + 1, v.begin() + 5); // v becomes \{1, 6\} Note: The
above methods do not change the capacity of the vector, only the size.
See Vector Size and Capacity. The erase method, which removes a range of
elements, is often used as a part of the erase-remove idiom. That is,
rst std::remove moves some elements to the end of the vector, and then
erase chops them o. This is a relatively inecient operation for any
indices less than the last index of the vector because all elements
after the erased segments must be relocated to new positions. For speed
critical applications that require ecient removal of arbitrary elements
in a container, see std::list. Deleting elements by value: std::vector
v\{ 1, 1, 2, 2, 3, 3 \}; int value\_to\_remove = 2;
v.erase(std::remove(v.begin(), v.end(), value\_to\_remove), v.end()); //
v becomes \{1, 1, 3, 3\} Deleting elements by condition: //
std::remove\_if needs a function, that takes a vector element as
argument and returns true, // if the element shall be removed bool
\_predicate(const int\& element) \{ return (element \textgreater{} 3);
// This will cause all elements to be deleted that are larger than 3 \}
std::vector v\{ 1, 2, 3, 4, 5, 6 \}; v.erase(std::remove\_if(v.begin(),
v.end(), \_predicate), v.end()); // v becomes \{1, 2, 3\} Deleting
elements by lambda, without creating additional predicate function
Version C++11 std::vector v\{ 1, 2, 3, 4, 5, 6 \};
v.erase(std::remove\_if(v.begin(), v.end(),
\href{auto\&\%20element}{}\{return element \textgreater{} 3;\} ),
v.end() ); Deleting elements by condition from a loop: std::vector v\{
1, 2, 3, 4, 5, 6 \}; std::vector::iterator it = v.begin(); while (it !=
v.end()) \{ if (condition) it = v.erase(it); // after erasing, `it' will
be set to the next element in v else ++it; // manually set `it' to the
next element in v \} While it is important not to increment it in case
of a deletion, you should consider using a dierent method when then
erasing repeatedly in a loop. Consider remove\_if for a more ecient way.
Deleting elements by condition from a reverse loop: std::vector v\{ -1,
0, 1, 2, 3, 4, 5, 6 \}; typedef std::vector::reverse\_iterator rev\_itr;
rev\_itr it = v.rbegin(); while (it != v.rend()) \{ // after the loop
only `0' will be in v int value = \emph{it; if (value) \{ ++it;  // See
explanation below for the following line. it =
rev\_itr(v.erase(it.base())); \} else ++it; \} Note some points for the
preceding loop: Given a reverse iterator it pointing to some element,
the method base gives the regular (non-reverse) iterator pointing to the
same element. vector::erase(iterator) erases the element pointed to by
an iterator, and returns an iterator to the element that followed the
given element. reverse\_iterator::reverse\_iterator(iterator) constructs
a reverse iterator from an iterator. Put altogether, the line it =
rev\_itr(v.erase(it.base())) says: take the reverse iterator it, have v
erase the element pointed by its regular iterator; take the resulting
iterator, construct a reverse iterator from it, and assign it to the
reverse iterator it. Deleting all elements using v.clear() does not free
up memory (capacity() of the vector remains unchanged). To reclaim
space, use: std::vector().swap(v); Version C++11 shrink\_to\_fit() frees
up unused vector capacity: v.shrink\_to\_fit(); The shrink\_to\_fit does
not guarantee to really reclaim space, but most current implementations
do. Section 49.4: Iterating Over std::vector You can iterate over a
std::vector in several ways. For each of the following sections, v is
dened as follows: std::vector v; Iterating in the Forward Direction
Version C++11 // Range based for for(const auto\& value: v) \{ std::cout
\textless\textless{} value \textless\textless{} ``\n''; \} // Using a
for loop with iterator for(auto it = std::begin(v); it != std::end(v);
++it) \{ std::cout \textless\textless{} }it \textless\textless{} ``\n'';
\} // Using for\_each algorithm, using a function or functor: void
fun(int const\& value) \{ std::cout \textless\textless{} value
\textless\textless{} ``\n''; \} std::for\_each(std::begin(v),
std::end(v), fun); // Using for\_each algorithm. Using a lambda:
std::for\_each(std::begin(v), std::end(v),
\href{int\%20const\&\%20value}{} \{ std::cout \textless\textless{} value
\textless\textless{} ``\n''; \}); Version \textless{} C++11 // Using a
for loop with iterator for(std::vector::iterator it = std::begin(v); it
!= std::end(v); ++it) \{ std::cout \textless\textless{} \emph{it
\textless\textless{} ``\n''; \} // Using a for loop with index
for(std::size\_t i = 0; i \textless{} v.size(); ++i) \{ std::cout
\textless\textless{} v{[}i{]} \textless\textless{} ``\n''; \} Iterating
in the Reverse Direction Version C++14 // There is no standard way to
use range based for for this. // See below for alternatives. // Using
for\_each algorithm // Note: Using a lambda for clarity. But a function
or functor will work std::for\_each(std::rbegin(v), std::rend(v),
\href{auto\%20const\&\%20value}{} \{ std::cout \textless\textless{}
value \textless\textless{} ``\n''; \}); // Using a for loop with
iterator for(auto rit = std::rbegin(v); rit != std::rend(v); ++rit) \{
std::cout \textless\textless{} }rit \textless\textless{} ``\n''; \} //
Using a for loop with index for(std::size\_t i = 0; i \textless{}
v.size(); ++i) \{ std::cout \textless\textless{} v{[}v.size() - 1 - i{]}
\textless\textless{} ``\n''; \} Though there is no built-in way to use
the range based for to reverse iterate; it is relatively simple to x
this. The range based for uses begin() and end() to get iterators and
thus simulating this with a wrapper object can achieve the results we
require. Version C++14 template struct ReverseRange \{ C c; // could be
a reference or a copy, if the original was a temporary
ReverseRange(C\&\& cin): c(std::forward(cin)) \{\}
ReverseRange(ReverseRange\&\&)=default; ReverseRange\&
operator=(ReverseRange\&\&)=delete; auto begin() const \{return
std::rbegin(c);\} auto end() const \{return std::rend(c);\} \}; // C is
meant to be deduced, and perfect forwarded into template ReverseRange
make\_ReverseRange(C\&\& c) \{return \{std::forward(c)\};\} int main()
\{ std::vector v \{ 1,2,3,4\}; for(auto const\& value:
make\_ReverseRange(v)) \{ std::cout \textless\textless{} value
\textless\textless{} ``\n''; \} \} Enforcing const elements Since C++11
the cbegin() and cend() methods allow you to obtain a constant iterator
for a vector, even if the vector is non-const. A constant iterator
allows you to read but not modify the contents of the vector which is
useful to enforce const correctness: Version C++11 // forward iteration
for (auto pos = v.cbegin(); pos != v.cend(); ++pos) \{ // type of pos is
vector::const\_iterator // \emph{pos = 5; // Compile error - can't write
via const iterator \} // reverse iteration for (auto pos = v.crbegin();
pos != v.crend(); ++pos) \{ // type of pos is vector::const\_iterator //
}pos = 5; // Compile error - can't write via const iterator \} //
expects Functor::operand()(T\&) for\_each(v.begin(), v.end(),
Functor()); // expects Functor::operand()(const T\&)
for\_each(v.cbegin(), v.cend(), Functor()) Version C++17 as\_const
extends this to range iteration: for (auto const\& e :
std::as\_const(v)) \{ std::cout \textless\textless{} e
\textless\textless{} `\n'; \} This is easy to implement in earlier
versions of C++: Version C++14 template constexpr std::add\_const\_t\&
as\_const(T\& t) noexcept \{ return t; \} A Note on Eciency Since the
class std::vector is basically a class that manages a dynamically
allocated contiguous array, the same principle explained here applies to
C++ vectors. Accessing the vector's content by index is much more ecient
when following the row-major order principle. Of course, each access to
the vector also puts its management content into the cache as well, but
as has been debated many times (notably here and here), the dierence in
performance for iterating over a std::vector compared to a raw array is
negligible. So the same principle of eciency for raw arrays in C also
applies for C++'s std::vector. Section 49.5: vector: The Exception To So
Many, So Many Rules The standard (section 23.3.7) species that a
specialization of vector is provided, which optimizes space by packing
the bool values, so that each takes up only one bit. Since bits aren't
addressable in C++, this means that several requirements on vector are
not placed on vector: The data stored is not required to be contiguous,
so a vector can't be passed to a C API which expects a bool array. at(),
operator {[}{]}, and dereferencing of iterators do not return a
reference to bool. Rather they return a proxy object that (imperfectly)
simulates a reference to a bool by overloading its assignment operators.
As an example, the following code may not be valid for std::vector,
because dereferencing an iterator does not return a reference: Version
C++11 std::vector v = \{true, false\}; for (auto \&b: v) \{ \} // error
Similarly, functions expecting a bool\& argument cannot be used with the
result of operator {[}{]} or at() applied to vector, or with the result
of dereferencing its iterator: void f(bool\& b); f(v{[}0{]}); // error
f(\emph{v.begin()); // error The implementation of std::vector is
dependent on both the compiler and architecture. The specialisation is
implemented by packing n Booleans into the lowest addressable section of
memory. Here, n is the size in bits of the lowest addressable memory. In
most modern systems this is 1 byte or 8 bits. This means that one byte
can store 8 Boolean values. This is an improvement over the traditional
implementation where 1 Boolean value is stored in 1 byte of memory.
Note: The below example shows possible bitwise values of individual
bytes in a traditional vs.~optimized vector. This will not always hold
true in all architectures. It is, however, a good way of visualising the
optimization. In the below examples a byte is represented as {[}x, x, x,
x, x, x, x, x{]}. Traditional std::vector storing 8 Boolean values:
Version C++11 std::vector trad\_vect = \{true, false, false, false,
true, false, true, true\}; Bitwise representation:
{[}0,0,0,0,0,0,0,1{]}, {[}0,0,0,0,0,0,0,0{]}, {[}0,0,0,0,0,0,0,0{]},
{[}0,0,0,0,0,0,0,0{]}, {[}0,0,0,0,0,0,0,1{]}, {[}0,0,0,0,0,0,0,0{]},
{[}0,0,0,0,0,0,0,1{]}, {[}0,0,0,0,0,0,0,1{]} Specialized std::vector
storing 8 Boolean values: Version C++11 std::vector optimized\_vect =
\{true, false, false, false, true, false, true, true\}; Bitwise
representation: {[}1,0,0,0,1,0,1,1{]} Notice in the above example, that
in the traditional version of std::vector, 8 Boolean values take up 8
bytes of memory, whereas in the optimized version of std::vector, they
only use 1 byte of memory. This is a signicant improvement on memory
usage. If you need to pass a vector to an C-style API, you may need to
copy the values to an array, or nd a better way to use the API, if
memory and performance are at risk. Section 49.6: Inserting Elements
Appending an element at the end of a vector (by copying/moving): struct
Point \{ double x, y;  Point(double x, double y) : x(x), y(y) \{\} \};
std::vector v; Point p(10.0, 2.0); v.push\_back(p); // p is copied into
the vector. Version C++11 Appending an element at the end of a vector by
constructing the element in place: std::vector v; v.emplace\_back(10.0,
2.0); // The arguments are passed to the constructor of the // given
type (here Point). The object is constructed // in the vector, avoiding
a copy. Note that std::vector does not have a push\_front() member
function due to performance reasons. Adding an element at the beginning
causes all existing elements in the vector to be moved. If you want to
frequently insert elements at the beginning of your container, then you
might want to use std::list or std::deque instead. Inserting an element
at any position of a vector: std::vector v\{ 1, 2, 3 \};
v.insert(v.begin(), 9); // v now contains \{9, 1, 2, 3\} Version C++11
Inserting an element at any position of a vector by constructing the
element in place: std::vector v\{ 1, 2, 3 \}; v.emplace(v.begin()+1, 9);
// v now contains \{1, 9, 2, 3\} Inserting another vector at any
position of the vector: std::vector v(4); // contains: 0, 0, 0, 0
std::vector v2(2, 10); // contains: 10, 10 v.insert(v.begin()+2,
v2.begin(), v2.end()); // contains: 0, 0, 10, 10, 0, 0 Inserting an
array at any position of a vector: std::vector v(4); // contains: 0, 0,
0, 0 int a {[}{]} = \{1, 2, 3\}; // contains: 1, 2, 3
v.insert(v.begin()+1, a, a+sizeof(a)/sizeof(a{[}0{]})); // contains: 0,
1, 2, 3, 0, 0, 0 Use reserve() before inserting multiple elements if
resulting vector size is known beforehand to avoid multiple
reallocations (see vector size and capacity): std::vector v;
v.reserve(100); for(int i = 0; i \textless{} 100; ++i)
v.emplace\_back(i); Be sure to not make the mistake of calling resize()
in this case, or you will inadvertently create a vector with 200
elements where only the latter one hundred will have the value you
intended. Section 49.7: Using std::vector as a C array There are several
ways to use a std::vector as a C array (for example, for compatibility
with C libraries). This is possible because the elements in a vector are
stored contiguously. Version C++11 std::vector v\{ 1, 2, 3 \}; int} p =
v.data(); In contrast to solutions based on previous C++ standards (see
below), the member function .data() may also be applied to empty
vectors, because it doesn't cause undened behavior in this case. Before
C++11, you would take the address of the vector's rst element to get an
equivalent pointer, if the vector isn't empty, these both methods are
interchangeable: int* p = \&v{[}0{]}; // combine subscript operator and
0 literal int* p = \&v.front(); // explicitly reference the first
element Note: If the vector is empty, v{[}0{]} and v.front() are undened
and cannot be used. When storing the base address of the vector's data,
note that many operations (such as push\_back, resize, etc.) can change
the data memory location of the vector, thus invalidating previous data
pointers. For example: std::vector v; int* p = v.data(); v.resize(42);
// internal memory location changed; value of p is now invalid Section
49.8: Finding an Element in std::vector The function std::find, dened in
the header, can be used to nd an element in a std::vector. std::find
uses the operator== to compare elements for equality. It returns an
iterator to the rst element in the range that compares equal to the
value. If the element in question is not found, std::find returns
std::vector::end (or std::vector::cend if the vector is const). Version
\textless{} C++11 static const int arr{[}{]} = \{5, 4, 3, 2, 1\};
std::vector v (arr, arr + sizeof(arr) / sizeof(arr{[}0{]}) );
std::vector::iterator it = std::find(v.begin(), v.end(), 4);
std::vector::difference\_type index = std::distance(v.begin(), it); //
\passthrough{\lstinline!it!} points to the second element of the vector,
\passthrough{\lstinline!index!} is 1 std::vector::iterator missing =
std::find(v.begin(), v.end(), 10); std::vector::difference\_type
index\_missing = std::distance(v.begin(), missing); //
\passthrough{\lstinline!missing!} is v.end(),
\passthrough{\lstinline!index\_missing!} is 5 (ie. size of the vector)
Version C++11 std::vector v \{ 5, 4, 3, 2, 1 \}; auto it =
std::find(v.begin(), v.end(), 4); auto index = std::distance(v.begin(),
it); // \passthrough{\lstinline!it!} points to the second element of the
vector, \passthrough{\lstinline!index!} is 1 auto missing =
std::find(v.begin(), v.end(), 10); auto index\_missing =
std::distance(v.begin(), missing); // \passthrough{\lstinline!missing!}
is v.end(), \passthrough{\lstinline!index\_missing!} is 5 (ie. size of
the vector) If you need to perform many searches in a large vector, then
you may want to consider sorting the vector rst, before using the
binary\_search algorithm. To nd the rst element in a vector that satises
a condition, std::find\_if can be used. In addition to the two
parameters given to std::find, std::find\_if accepts a third argument
which is a function object or function pointer to a predicate function.
The predicate should accept an element from the container as an argument
and return a value convertible to bool, without modifying the container:
Version \textless{} C++11 bool isEven(int val) \{ return (val \% 2 ==
0); \} struct moreThan \{ moreThan(int limit) : \_limit(limit) \{\} bool
operator()(int val) \{ return val \textgreater{} \_limit; \} int
\_limit; \}; static const int arr{[}{]} = \{1, 3, 7, 8\}; std::vector v
(arr, arr + sizeof(arr) / sizeof(arr{[}0{]}) ); std::vector::iterator it
= std::find\_if(v.begin(), v.end(), isEven); //
\passthrough{\lstinline!it!} points to 8, the first even element
std::vector::iterator missing = std::find\_if(v.begin(), v.end(),
moreThan(10)); // \passthrough{\lstinline!missing!} is v.end(), as no
element is greater than 10 Version C++11 // find the first value that is
even std::vector v = \{1, 3, 7, 8\}; auto it = std::find\_if(v.begin(),
v.end(), \href{int\%20val}{}\{return val \% 2 == 0;\}); //
\passthrough{\lstinline!it!} points to 8, the first even element auto
missing = std::find\_if(v.begin(), v.end(), \href{int\%20val}{}\{return
val \textgreater{} 10;\}); // \passthrough{\lstinline!missing!} is
v.end(), as no element is greater than 10 Section 49.9: Concatenating
Vectors One std::vector can be append to another by using the member
function insert(): std::vector a = \{0, 1, 2, 3, 4\}; std::vector b =
\{5, 6, 7, 8, 9\}; a.insert(a.end(), b.begin(), b.end()); However, this
solution fails if you try to append a vector to itself, because the
standard species that iterators given to insert() must not be from the
same range as the receiver object's elements. Version c++11 Instead of
using the vector's member functions, the functions std::begin() and
std::end() can be used: a.insert(std::end(a), std::begin(b),
std::end(b)); This is a more general solution, for example, because b
can also be an array. However, also this solution doesn't allow you to
append a vector to itself. If the order of the elements in the receiving
vector doesn't matter, considering the number of elements in each vector
can avoid unnecessary copy operations: if (b.size() \textless{}
a.size()) a.insert(a.end(), b.begin(), b.end()); else b.insert(b.end(),
a.begin(), a.end()); Section 49.10: Matrices Using Vectors Vectors can
be used as a 2D matrix by dening them as a vector of vectors. A matrix
with 3 rows and 4 columns with each cell initialised as 0 can be dened
as: std::vector\textless std::vector \textgreater{} matrix(3,
std::vector(4)); Version C++11 The syntax for initializing them using
initialiser lists or otherwise are similar to that of a normal vector.
std::vector\textless std::vector\textgreater{} matrix = \{ \{0,1,2,3\},
\{4,5,6,7\}, \{8,9,10,11\} \}; Values in such a vector can be accessed
similar to a 2D array int var = matrix{[}0{]}{[}2{]}; Iterating over the
entire matrix is similar to that of a normal vector but with an extra
dimension. for(int i = 0; i \textless{} 3; ++i) \{ for(int j = 0; j
\textless{} 4; ++j) \{ std::cout \textless\textless{}
matrix{[}i{]}{[}j{]} \textless\textless{} std::endl; \} \} Version C++11
for(auto\& row: matrix) \{ for(auto\& col : row) \{ std::cout
\textless\textless{} col \textless\textless{} std::endl; \} \} A vector
of vectors is a convenient way to represent a matrix but it's not the
most ecient: individual vectors are scattered around memory and the data
structure isn't cache friendly. Also, in a proper matrix, the length of
every row must be the same (this isn't the case for a vector of
vectors). The additional exibility can be a source of errors. Section
49.11: Using a Sorted Vector for Fast Element Lookup The header provides
a number of useful functions for working with sorted vectors. An
important prerequisite for working with sorted vectors is that the
stored values are comparable with \textless. An unsorted vector can be
sorted by using the function std::sort(): std::vector v; // add some
code here to fill v with some elements std::sort(v.begin(), v.end());
Sorted vectors allow ecient element lookup using the function
std::lower\_bound(). Unlike std::find(), this performs an ecient binary
search on the vector. The downside is that it only gives valid results
for sorted input ranges: // search the vector for the first element with
value 42 std::vector::iterator it = std::lower\_bound(v.begin(),
v.end(), 42); if (it != v.end() \&\& *it == 42) \{ // we found the
element! \} Note: If the requested value is not part of the vector,
std::lower\_bound() will return an iterator to the rst element that is
greater than the requested value. This behavior allows us to insert a
new element at its right place in an already sorted vector: int const
new\_element = 33; v.insert(std::lower\_bound(v.begin(), v.end(),
new\_element), new\_element); If you need to insert a lot of elements at
once, it might be more ecient to call push\_back() for all them rst and
then call std::sort() once all elements have been inserted. In this
case, the increased cost of the sorting can pay o against the reduced
cost of inserting new elements at the end of the vector and not in the
middle. If your vector contains multiple elements of the same value,
std::lower\_bound() will try to return an iterator to the rst element of
the searched value. However, if you need to insert a new element after
the last element of the searched value, you should use the function
std::upper\_bound() as this will cause less shifting around of elements:
v.insert(std::upper\_bound(v.begin(), v.end(), new\_element),
new\_element); If you need both the upper bound and the lower bound
iterators, you can use the function std::equal\_range() to retrieve both
of them eciently with one call:
std::pair\textless std::vector::iterator,
std::vector::iterator\textgreater{} rg = std::equal\_range(v.begin(),
v.end(), 42); std::vector::iterator lower\_bound = rg.first;
std::vector::iterator upper\_bound = rg.second; In order to test whether
an element exists in a sorted vector (although not specic to vectors),
you can use the function std::binary\_search(): bool exists =
std::binary\_search(v.begin(), v.end(), value\_to\_find); Section
49.12: Reducing the Capacity of a Vector A std::vector automatically
increases its capacity upon insertion as needed, but it never reduces
its capacity after element removal. // Initialize a vector with 100
elements std::vector v(100); // The vector's capacity is always at least
as large as its size auto const old\_capacity = v.capacity(); //
old\_capacity \textgreater= 100 // Remove half of the elements
v.erase(v.begin() + 50, v.end()); // Reduces the size from 100 to 50
(v.size() == 50), // but not the capacity (v.capacity() ==
old\_capacity) To reduce its capacity, we can copy the contents of a
vector to a new temporary vector. The new vector will have the minimum
capacity that is needed to store all elements of the original vector. If
the size reduction of the original vector was signicant, then the
capacity reduction for the new vector is likely to be signicant. We can
then swap the original vector with the temporary one to retain its
minimized capacity: std::vector(v).swap(v); Version C++11 In C++11 we
can use the shrink\_to\_fit() member function for a similar eect:
v.shrink\_to\_fit(); Note: The shrink\_to\_fit() member function is a
request and doesn't guarantee to reduce capacity. Section 49.13: Vector
size and capacity Vector size is simply the number of elements in the
vector: 1. Current vector size is queried by size() member function.
Convenience empty() function returns true if size is 0: vector v = \{ 1,
2, 3 \}; // size is 3 const vector::size\_type size = v.size(); cout
\textless\textless{} size \textless\textless{} endl; // prints 3 cout
\textless\textless{} boolalpha \textless\textless{} v.empty()
\textless\textless{} endl; // prints false 2. Default constructed vector
starts with a size of 0: vector v; // size is 0 cout
\textless\textless{} v.size() \textless\textless{} endl; // prints 0 3.
Adding N elements to vector increases size by N (e.g.~by push\_back(),
insert() or resize() functions). 4. Removing N elements from vector
decreases size by N (e.g.~by pop\_back(), erase() or clear() functions).
5. Vector has an implementation-specic upper limit on its size, but you
are likely to run out of RAM before reaching it: vector v; const
vector::size\_type max\_size = v.max\_size(); cout \textless\textless{}
max\_size \textless\textless{} endl; // prints some large number
v.resize( max\_size ); // probably won't work v.push\_back( 1 ); //
definitely won't work Common mistake: size is not necessarily (or even
usually) int: // !!!bad!!!evil!!! vector v\_bad( N, 1 ); // constructs
large N size vector for( int i = 0; i \textless{} v\_bad.size(); ++i )
\{ // size is not supposed to be int! do\_something( v\_bad{[}i{]} ); \}
Vector capacity diers from size. While size is simply how many elements
the vector currently has, capacity is for how many elements it
allocated/reserved memory for. That is useful, because too frequent
(re)allocations of too large sizes can be expensive. 1. Current vector
capacity is queried by capacity() member function. Capacity is always
greater or equal to size: vector v = \{ 1, 2, 3 \}; // size is 3,
capacity is \textgreater= 3 const vector::size\_type capacity =
v.capacity(); cout \textless\textless{} capacity \textless\textless{}
endl; // prints number \textgreater= 3 2. You can manually reserve
capacity by reserve( N ) function (it changes vector capacity to N): //
!!!bad!!!evil!!! vector v\_bad; for( int i = 0; i \textless{} 10000; ++i
) \{ v\_bad.push\_back( i ); // possibly lot of reallocations \} // good
vector v\_good; v\_good.reserve( 10000 ); // good! only one allocation
for( int i = 0; i \textless{} 10000; ++i ) \{ v\_good.push\_back( i );
// no allocations needed anymore \} 3. You can request for the excess
capacity to be released by shrink\_to\_fit() (but the implementation
doesn't have to obey you). This is useful to conserve used memory:
vector v = \{ 1, 2, 3, 4, 5 \}; // size is 5, assume capacity is 6
v.shrink\_to\_fit(); // capacity is 5 (or possibly still 6) cout
\textless\textless{} boolalpha \textless\textless{} v.capacity() ==
v.size() \textless\textless{} endl; // prints likely true (but possibly
false) Vector partly manages capacity automatically, when you add
elements it may decide to grow. Implementers like to use 2 or 1.5 for
the grow factor (golden ratio would be the ideal value - but is
impractical due to being rational number). On the other hand vector
usually do not automatically shrink. For example: vector v; // capacity
is possibly (but not guaranteed) to be 0 v.push\_back( 1 ); // capacity
is some starter value, likely 1 v.clear(); // size is 0 but capacity is
still same as before!

\hypertarget{professional-notes-chapter-50-stdmap-1}{%
\chapter{Professional Notes: Chapter 50:
std::map}\label{professional-notes-chapter-50-stdmap-1}}

To use any of std::map or std::multimap the header le should be
included. std::map and std::multimap both keep their elements sorted
according to the ascending order of keys. In case of std::multimap, no
sorting occurs for the values of the same key. The basic dierence
between std::map and std::multimap is that the std::map one does not
allow duplicate values for the same key where std::multimap does. Maps
are implemented as binary search trees. So search(), insert(), erase()
takes (log n) time in average. For constant time operation use
std::unordered\_map. size() and empty() functions have (1) time
complexity, number of nodes is cached to avoid walking through tree each
time these functions are called. Section 50.1: Accessing elements An
std::map takes (key, value) pairs as input. Consider the following
example of std::map initialization: std::map \textless{} std::string,
int \textgreater{} ranking \{ std::make\_pair(``stackoverflow'', 2),
std::make\_pair(``docs-beta'', 1) \}; In an std::map , elements can be
inserted as follows: ranking{[}``stackoverflow''{]}=2;
ranking{[}``docs-beta''{]}=1; In the above example, if the key
stackoverflow is already present, its value will be updated to 2. If it
isn't already present, a new entry will be created. In an std::map,
elements can be accessed directly by giving the key as an index:
std::cout \textless\textless{} ranking{[} ``stackoverflow'' {]}
\textless\textless{} std::endl; Note that using the operator{[}{]} on
the map will actually insert a new value with the queried key into the
map. This means that you cannot use it on a const std::map, even if the
key is already stored in the map. To prevent this insertion, check if
the element exists (for example by using find()) or use at() as
described below. Version C++11 Elements of a std::map can be accessed
with at(): std::cout \textless\textless{} ranking.at(``stackoverflow'')
\textless\textless{} std::endl; Note that at() will throw an
std::out\_of\_range exception if the container does not contain the
requested element. In both containers std::map and std::multimap,
elements can be accessed using iterators: Version C++11 // Example using
begin() std::multimap \textless{} int, std::string \textgreater{} mmp
\{ std::make\_pair(2, ``stackoverflow''), std::make\_pair(1,
``docs-beta''), std::make\_pair(2, ``stackexchange'') \}; auto it =
mmp.begin(); std::cout \textless\textless{} it-\textgreater first
\textless\textless{} " : " \textless\textless{} it-\textgreater second
\textless\textless{} std::endl; // Output: ``1 : docs-beta'' it++;
std::cout \textless\textless{} it-\textgreater first
\textless\textless{} " : " \textless\textless{} it-\textgreater second
\textless\textless{} std::endl; // Output: ``2 : stackoverflow'' it++;
std::cout \textless\textless{} it-\textgreater first
\textless\textless{} " : " \textless\textless{} it-\textgreater second
\textless\textless{} std::endl; // Output: ``2 : stackexchange'' //
Example using rbegin() std::map \textless{} int, std::string
\textgreater{} mp \{ std::make\_pair(2, ``stackoverflow''),
std::make\_pair(1, ``docs-beta''), std::make\_pair(2, ``stackexchange'')
\}; auto it2 = mp.rbegin(); std::cout \textless\textless{}
it2-\textgreater first \textless\textless{} " : " \textless\textless{}
it2-\textgreater second \textless\textless{} std::endl; // Output: ``2 :
stackoverflow'' it2++; std::cout \textless\textless{}
it2-\textgreater first \textless\textless{} " : " \textless\textless{}
it2-\textgreater second \textless\textless{} std::endl; // Output: ``1 :
docs-beta'' Section 50.2: Inserting elements An element can be inserted
into a std::map only if its key is not already present in the map. Given
for example: std::map\textless{} std::string, size\_t \textgreater{}
fruits\_count; A key-value pair is inserted into a std::map through the
insert() member function. It requires a pair as an argument:
fruits\_count.insert(\{``grapes'', 20\});
fruits\_count.insert(make\_pair(``orange'', 30));
fruits\_count.insert(pair\textless std::string,
size\_t\textgreater(``banana'', 40));
fruits\_count.insert(map\textless std::string,
size\_t\textgreater::value\_type(``cherry'', 50)); The insert() function
returns a pair consisting of an iterator and a bool value: If the
insertion was successful, the iterator points to the newly inserted
element, and the bool value is true. If there was already an element
with the same key, the insertion fails. When that happens, the iterator
points to the element causing the conict, and the bool is value is
false. The following method can be used to combine insertion and
searching operation: auto success = fruits\_count.insert(\{``grapes'',
20\}); if (!success.second) \{ // we already have `grapes' in the map
success.first-\textgreater second += 20; // access the iterator to
update the value \} For convenience, the std::map container provides the
subscript operator to access elements in the map and to insert new ones
if they don't exist: fruits\_count{[}``apple''{]} = 10; While simpler,
it prevents the user from checking if the element already exists. If an
element is missing, std::map::operator{[}{]} implicitly creates it,
initializing it with the default constructor before overwriting it with
the supplied value. insert() can be used to add several elements at
once using a braced list of pairs. This version of insert() returns
void: fruits\_count.insert(\{\{``apricot'', 1\}, \{``jackfruit'', 1\},
\{``lime'', 1\}, \{``mango'', 7\}\}); insert() can also be used to add
elements by using iterators denoting the begin and end of value\_type
values: std::map\textless{} std::string, size\_t \textgreater{}
fruit\_list\{ \{``lemon'', 0\}, \{``olive'', 0\}, \{``plum'', 0\}\};
fruits\_count.insert(fruit\_list.begin(), fruit\_list.end()); Example:
std::map\textless std::string, size\_t\textgreater{} fruits\_count;
std::string fruit; while(std::cin \textgreater\textgreater{} fruit)\{ //
insert an element with `fruit' as key and `1' as value // (if the key is
already stored in fruits\_count, insert does nothing) auto ret =
fruits\_count.insert(\{fruit, 1\}); if(!ret.second)\{ // `fruit' is
already in the map ++ret.first-\textgreater second; // increment the
counter \} \} Time complexity for an insertion operation is O(log n)
because std::map are implemented as trees. Version C++11 A pair can be
constructed explicitly using make\_pair() and emplace():
std::map\textless{} std::string , int \textgreater{} runs;
runs.emplace(``Babe Ruth'', 714); runs.insert(make\_pair(``Barry
Bonds'', 762)); If we know where the new element will be inserted, then
we can use emplace\_hint() to specify an iterator hint. If the new
element can be inserted just before hint, then the insertion can be done
in constant time. Otherwise it behaves in the same way as emplace():
std::map\textless{} std::string , int \textgreater{} runs; auto it =
runs.emplace(``Barry Bonds'', 762); // get iterator to the inserted
element // the next element will be before ``Barry Bonds'', so it is
inserted before `it' runs.emplace\_hint(it, ``Babe Ruth'', 714); Section
50.3: Searching in std::map or in std::multimap There are several ways
to search a key in std::map or in std::multimap. To get the iterator of
the rst occurrence of a key, the find() function can be used. It returns
end() if the key does not exist. std::multimap\textless{} int , int
\textgreater{} mmp\{ \{1, 2\}, \{3, 4\}, \{6, 5\}, \{8, 9\}, \{3, 4\},
\{6, 7\} \}; auto it = mmp.find(6); if(it!=mmp.end()) std::cout
\textless\textless{} it-\textgreater first \textless\textless{} ``,''
\textless\textless{} it-\textgreater second \textless\textless{}
std::endl; //prints: 6, 5 else  std::cout \textless\textless{} ``Value
does not exist!'' \textless\textless{} std::endl; it = mmp.find(66);
if(it!=mmp.end()) std::cout \textless\textless{} it-\textgreater first
\textless\textless{} ``,'' \textless\textless{} it-\textgreater second
\textless\textless{} std::endl; else std::cout \textless\textless{}
``Value does not exist!'' \textless\textless{} std::endl; // This line
would be executed. Another way to nd whether an entry exists in std::map
or in std::multimap is using the count() function, which counts how many
values are associated with a given key. Since std::map associates only
one value with each key, its count() function can only return 0 (if the
key is not present) or 1 (if it is). For std::multimap, count() can
return values greater than 1 since there can be several values
associated with the same key. std::map\textless{} int , int
\textgreater{} mp\{ \{1, 2\}, \{3, 4\}, \{6, 5\}, \{8, 9\}, \{3, 4\},
\{6, 7\} \}; if(mp.count(3) \textgreater{} 0) // 3 exists as a key in
map std::cout \textless\textless{} ``The key exists!''
\textless\textless{} std::endl; // This line would be executed. else
std::cout \textless\textless{} ``The key does not exist!''
\textless\textless{} std::endl; If you only care whether some element
exists, find is strictly better: it documents your intent and, for
multimaps, it can stop once the rst matching element has been found. In
the case of std::multimap, there could be several elements having the
same key. To get this range, the equal\_range() function is used which
returns std::pair having iterator lower bound (inclusive) and upper
bound (exclusive) respectively. If the key does not exist, both
iterators would point to end(). auto eqr = mmp.equal\_range(6); auto st
= eqr.first, en = eqr.second; for(auto it = st; it != en; ++it)\{
std::cout \textless\textless{} it-\textgreater first
\textless\textless{} ``,'' \textless\textless{} it-\textgreater second
\textless\textless{} std::endl; \} // prints: 6, 5 // 6, 7 Section 50.4:
Initializing a std::map or std::multimap std::map and std::multimap both
can be initialized by providing key-value pairs separated by comma.
Key-value pairs could be provided by either \{key, value\} or can be
explicitly created by std::make\_pair(key, value). As std::map does not
allow duplicate keys and comma operator performs right to left, the pair
on right would be overwritten with the pair with same key on the left.
std::multimap \textless{} int, std::string \textgreater{} mmp \{
std::make\_pair(2, ``stackoverflow''), std::make\_pair(1,
``docs-beta''), std::make\_pair(2, ``stackexchange'') \}; // 1 docs-beta
// 2 stackoverflow // 2 stackexchange std::map \textless{} int,
std::string \textgreater{} mp \{ std::make\_pair(2, ``stackoverflow''),
std::make\_pair(1, ``docs-beta''), std::make\_pair(2, ``stackexchange'')
\}; // 1 docs-beta // 2 stackoverflow Both could be initialized with
iterator. // From std::map or std::multimap iterator
std::multimap\textless{} int , int \textgreater{} mmp\{ \{1, 2\}, \{3,
4\}, \{6, 5\}, \{8, 9\}, \{6, 8\}, \{3, 4\}, \{6, 7\} \}; // \{1, 2\},
\{3, 4\}, \{3, 4\}, \{6, 5\}, \{6, 8\}, \{6, 7\}, \{8, 9\} auto it =
mmp.begin(); std::advance(it,3); //moved cursor on first \{6, 5\}
std::map\textless{} int, int \textgreater{} mp(it, mmp.end()); // \{6,
5\}, \{8, 9\} //From std::pair array std::pair\textless{} int, int
\textgreater{} arr{[}10{]}; arr{[}0{]} = \{1, 3\}; arr{[}1{]} = \{1,
5\}; arr{[}2{]} = \{2, 5\}; arr{[}3{]} = \{0, 1\}; std::map\textless{}
int, int \textgreater{} mp(arr,arr+4); //\{0 , 1\}, \{1, 3\}, \{2, 5\}
//From std::vector of std::pair std::vector\textless{}
std::pair\textless int, int\textgreater{} \textgreater{} v\{ \{1, 5\},
\{5, 1\}, \{3, 6\}, \{3, 2\} \}; std::multimap\textless{} int, int
\textgreater{} mp(v.begin(), v.end()); // \{1, 5\}, \{3, 6\}, \{3, 2\},
\{5, 1\} Section 50.5: Checking number of elements The container
std::map has a member function empty(), which returns true or false,
depending on whether the map is empty or not. The member function size()
returns the number of element stored in a std::map container:
std::map\textless std::string , int\textgreater{} rank
\{\{``facebook.com'', 1\} ,\{``google.com'', 2\}, \{``youtube.com'',
3\}\}; if(!rank.empty())\{ std::cout \textless\textless{} ``Number of
elements in the rank map:'' \textless\textless{} rank.size()
\textless\textless{} std::endl; \} else\{ std::cout \textless\textless{}
``The rank map is empty'' \textless\textless{} std::endl; \} Section
50.6: Types of Maps Regular Map A map is an associative container,
containing key-value pairs. \#include \#include
std::map\textless std::string, size\_t\textgreater{} fruits\_count; In
the above example, std::string is the key type, and size\_t is a value.
The key acts as an index in the map. Each key must be unique, and must
be ordered. If you need mutliple elements with the same key, consider
using multimap (explained below) If your value type does not specify any
ordering, or you want to override the default ordering, you may provide
one: \#include \#include \#include struct StrLess \{ bool
operator()(const std::string\& a, const std::string\& b) \{ return
strncmp(a.c\_str(), b.c\_str(), 8)\textless0; //compare only up to 8
first characters \} \} std::map\textless std::string, size\_t,
StrLess\textgreater{} fruits\_count2; If StrLess comparator returns
false for two keys, they are considered the same even if their actual
contents dier. Multi-Map Multimap allows multiple key-value pairs with
the same key to be stored in the map. Otherwise, its interface and
creation is very similar to the regular map. \#include \#include
std::multimap\textless std::string, size\_t\textgreater{} fruits\_count;
std::multimap\textless std::string, size\_t, StrLess\textgreater{}
fruits\_count2; Hash-Map (Unordered Map) A hash map stores key-value
pairs similar to a regular map. It does not order the elements with
respect to the key though. Instead, a hash value for the key is used to
quickly access the needed key-value pairs. \#include \#include
std::unordered\_map\textless std::string, size\_t\textgreater{}
fruits\_count; Unordered maps are usually faster, but the elements are
not stored in any predictable order. For example, iterating over all
elements in an unordered\_map gives the elements in a seemingly random
order. Section 50.7: Deleting elements Removing all elements:
std::multimap\textless{} int , int \textgreater{} mmp\{ \{1, 2\}, \{3,
4\}, \{6, 5\}, \{8, 9\}, \{3, 4\}, \{6, 7\} \}; mmp.clear(); //empty
multimap Removing element from somewhere with the help of iterator:
std::multimap\textless{} int , int \textgreater{} mmp\{ \{1, 2\}, \{3,
4\}, \{6, 5\}, \{8, 9\}, \{3, 4\}, \{6, 7\} \}; // \{1, 2\}, \{3, 4\},
\{3, 4\}, \{6, 5\}, \{6, 7\}, \{8, 9\} auto it = mmp.begin();
std::advance(it,3); // moved cursor on first \{6, 5\} mmp.erase(it); //
\{1, 2\}, \{3, 4\}, \{3, 4\}, \{6, 7\}, \{8, 9\} Removing all elements
in a range: std::multimap\textless{} int , int \textgreater{} mmp\{ \{1,
2\}, \{3, 4\}, \{6, 5\}, \{8, 9\}, \{3, 4\}, \{6, 7\} \}; // \{1, 2\},
\{3, 4\}, \{3, 4\}, \{6, 5\}, \{6, 7\}, \{8, 9\} auto it = mmp.begin();
auto it2 = it; it++; //moved first cursor on first \{3, 4\}
std::advance(it2,3); //moved second cursor on first \{6, 5\}
mmp.erase(it,it2); // \{1, 2\}, \{6, 5\}, \{6, 7\}, \{8, 9\} Removing
all elements having a provided value as key: std::multimap\textless{}
int , int \textgreater{} mmp\{ \{1, 2\}, \{3, 4\}, \{6, 5\}, \{8, 9\},
\{3, 4\}, \{6, 7\} \}; // \{1, 2\}, \{3, 4\}, \{3, 4\}, \{6, 5\}, \{6,
7\}, \{8, 9\} mmp.erase(6); // \{1, 2\}, \{3, 4\}, \{3, 4\}, \{8, 9\}
Removing elements that satisfy a predicate pred:
std::map\textless int,int\textgreater{} m; auto it = m.begin(); while
(it != m.end()) \{ if (pred(*it)) it = m.erase(it); else ++it; \}
Section 50.8: Iterating over std::map or std::multimap std::map or
std::multimap could be traversed by the following ways:
std::multimap\textless{} int , int \textgreater{} mmp\{ \{1, 2\}, \{3,
4\}, \{6, 5\}, \{8, 9\}, \{3, 4\}, \{6, 7\} \}; //Range based loop -
since C++11 for(const auto \&x: mmp) std::cout\textless\textless{}
x.first \textless\textless{}``:''\textless\textless{} x.second
\textless\textless{} std::endl; //Forward iterator for loop: it would
loop through first element to last element //it will be a
std::map\textless{} int, int \textgreater::iterator for (auto it =
mmp.begin(); it != mmp.end(); ++it) std::cout\textless\textless{}
it-\textgreater first \textless\textless{}``:''\textless\textless{}
it-\textgreater second \textless\textless{} std::endl; //Do something
with iterator //Backward iterator for loop: it would loop through last
element to first element //it will be a std::map\textless{} int, int
\textgreater::reverse\_iterator for (auto it = mmp.rbegin(); it !=
mmp.rend(); ++it) std::cout\textless\textless{} it-\textgreater first
\textless\textless" ``\textless\textless{} it-\textgreater second
\textless\textless{} std::endl; //Do something with iterator While
iterating over a std::map or a std::multimap, the use of auto is
preferred to avoid useless implicit conversions (see this SO answer for
more details). Section 50.9: Creating std::map with user-dened types as
key In order to be able to use a class as the key in a map, all that is
required of the key is that it be copiable and assignable. The ordering
within the map is dened by the third argument to the template (and the
argument to the constructor, if used). This defaults to std::less, which
defaults to the \textless{} operator, but there's no requirement to use
the defaults. Just write a comparison operator (preferably as a
functional object): struct CmpMyType \{ bool operator()( MyType const\&
lhs, MyType const\& rhs ) const \{  // \ldots{} \} \}; In C++,
the''compare" predicate must be a strict weak ordering. In particular,
compare(X,X) must return false for any X. i.e.~if CmpMyType()(a, b)
returns true, then CmpMyType()(b, a) must return false, and if both
return false, the elements are considered equal (members of the same
equivalence class). Strict Weak Ordering This is a mathematical term to
dene a relationship between two objects. Its denition is: Two objects x
and y are equivalent if both f(x, y) and f(y, x) are false. Note that an
object is always (by the irreexivity invariant) equivalent to itself. In
terms of C++ this means if you have two objects of a given type, you
should return the following values when compared with the operator
\textless. X a; X b; Condition: Test: Result a is equivalent to b: a
\textless{} b false a is equivalent to b b \textless{} a false a is less
than b a \textless{} b true a is less than b b \textless{} a false b is
less than a a \textless{} b false b is less than a b \textless{} a true
How you dene equivalent/less is totally dependent on the type of your
object.

\hypertarget{professional-notes-chapter-62-standard-library-algorithms-1}{%
\chapter{Professional Notes: Chapter 62: Standard Library
Algorithms}\label{professional-notes-chapter-62-standard-library-algorithms-1}}

Section 62.1: std::next\_permutation template\textless{} class Iterator
\textgreater{} bool next\_permutation( Iterator first, Iterator last );
template\textless{} class Iterator, class Compare \textgreater{} bool
next\_permutation( Iterator first, Iterator last, Compare cmpFun );
Eects: Sift the data sequence of the range {[}rst, last) into the next
lexicographically higher permutation. If cmpFun is provided, the
permutation rule is customized. Parameters: first- the beginning of the
range to be permutated, inclusive last - the end of the range to be
permutated, exclusive Return Value: Returns true if such permutation
exists. Otherwise the range is swaped to the lexicographically smallest
permutation and return false. Complexity: O(n), n is the distance from
first to last. Example: std::vector\textless{} int \textgreater{} v \{
1, 2, 3 \}; do \{ for( int i = 0; i \textless{} v.size(); i += 1 ) \{
std::cout \textless\textless{} v{[}i{]}; \} std::cout
\textless\textless{} std::endl; \}while( std::next\_permutation(
v.begin(), v.end() ) ); print all the permutation cases of 1,2,3 in
lexicographically-increasing order. output: Section 62.2: std::for\_each
template\textless class InputIterator, class Function\textgreater{}
Function for\_each(InputIterator first, InputIterator last, Function f);
Eects: Applies f to the result of dereferencing every iterator in the
range {[}first, last) starting from first and proceeding to last - 1.
Parameters: first, last - the range to apply f to. f - callable object
which is applied to the result of dereferencing every iterator in the
range {[}first, last). Return value: f (until C++11) and std::move(f)
(since C++11). Complexity: Applies f exactly last - first times.
Example: Version c++11 std::vector v \{ 1, 2, 4, 8, 16 \};
std::for\_each(v.begin(), v.end(), \href{int\%20elem}{} \{ std::cout
\textless\textless{} elem \textless\textless{} " ``; \}); Applies the
given function for every element of the vector v printing this element
to stdout. Section 62.3: std::accumulate Dened in header
template\textless class InputIterator, class T\textgreater{} T
accumulate(InputIterator first, InputIterator last, T init); // (1)
template\textless class InputIterator, class T, class
BinaryOperation\textgreater{} T accumulate(InputIterator first,
InputIterator last, T init, BinaryOperation f); // (2) Eects:
std::accumulate performs fold operation using f function on range
{[}first, last) starting with init as accumulator value. Eectively it's
equivalent of: T acc = init; for (auto it = first; first != last; ++it)
acc = f(acc, \emph{it); return acc; In version (1) operator+ is used in
place of f, so accumulate over container is equivalent of sum of
container elements. Parameters: first, last - the range to apply f to.
init - initial value of accumulator. f - binary folding function. Return
value: Accumulated value of f applications. Complexity: O(nk), where n
is the distance from first to last, O(k) is complexity of f function.
Example: Simple sum example: std::vector v \{ 2, 3, 4 \}; auto sum =
std::accumulate(v.begin(), v.end(), 1); std::cout \textless\textless{}
sum \textless\textless{} std::endl; Output: Convert digits to number:
Version \textless{} c++11 class Converter \{ public: int operator()(int
a, int d) const \{ return a } 10 + d; \} \}; and later const int
ds{[}3{]} = \{1, 2, 3\}; int n = std::accumulate(ds, ds + 3, 0,
Converter()); std::cout \textless\textless{} n \textless\textless{}
std::endl; Version c++11 const std::vector ds = \{1, 2, 3\}; int n =
std::accumulate(ds.begin(), ds.end(), 0, \href{int\%20a,\%20int\%20d}{}
\{ return a * 10 + d; \}); std::cout \textless\textless{} n
\textless\textless{} std::endl; Output: Section 62.4: std::nd template
\textless class InputIterator, class T\textgreater{} InputIterator find
(InputIterator first, InputIterator last, const T\& val); Eects Finds
the rst occurrence of val within the range {[}rst, last) Parameters
first =\textgreater{} iterator pointing to the beginning of the range
last =\textgreater{} iterator pointing to the end of the range val
=\textgreater{} The value to nd within the range Return An iterator
that points to the rst element within the range that is equal(==) to
val, the iterator points to last if val is not found. Example \#include
\#include \#include using namespace std; int main(int argc, const char *
argv{[}{]}) \{ //create a vector vector intVec \{4, 6, 8, 9, 10, 30,
55,100, 45, 2, 4, 7, 9, 43, 48\}; //define iterators vector::iterator
itr\_9; vector::iterator itr\_43; vector::iterator itr\_50; //calling
find itr\_9 = find(intVec.begin(), intVec.end(), 9); //occurs twice
itr\_43 = find(intVec.begin(), intVec.end(), 43); //occurs once //a
value not in the vector itr\_50 = find(intVec.begin(), intVec.end(),
50); //does not occur cout \textless\textless{}''first occurrence of: "
\textless\textless{} \emph{itr\_9 \textless\textless{} endl; cout
\textless\textless{} ``only occurrence of:'' \textless\textless{}
}itr\_43 \textless\textless{} Lendl; /\emph{ let's prove that itr\_9 is
pointing to the first occurrence of 9 by looking at the element after 9,
which should be 10 not 43 }/ cout \textless\textless{} ``element after
first 9:'' \textless\textless{} \emph{(itr\_9 + 1) \textless\textless{}
ends; /} to avoid dereferencing intVec.end(), lets look at the element
right before the end \emph{/ cout \textless\textless{} ``last element:''
\textless\textless{} }(itr\_50 - 1) \textless\textless{} endl; return 0;
\} Output first occurrence of: 9 only occurrence of: 43 element after
first 9: 10 last element: 48 Section 62.5: std::min\_element template
ForwardIterator min\_element (ForwardIterator first, ForwardIterator
last); template \textless class ForwardIterator, class
Compare\textgreater{} ForwardIterator min\_element (ForwardIterator
first, ForwardIterator last,Compare comp); Eects Finds the minimum
element in a range Parameters first - iterator pointing to the beginning
of the range last - iterator pointing to the end of the range comp - a
function pointer or function object that takes two arguments and returns
true or false indicating whether argument is less than argument 2. This
function should not modify inputs Return Iterator to the minimum element
in the range Complexity Linear in one less than the number of elements
compared. Example \#include \#include \#include \#include //to use
make\_pair using namespace std; //function compare two pairs bool
pairLessThanFunction(const pair\textless string, int\textgreater{} \&p1,
const pair\textless string, int\textgreater{} \&p2) \{ return p1.second
\textless{} p2.second; \} int main(int argc, const char * argv{[}{]}) \{
vector intVec \{30,200,167,56,75,94,10,73,52,6,39,43\};
vector\textless pair\textless string, int\textgreater\textgreater{}
pairVector = \{make\_pair(``y'', 25), make\_pair(``b'', 2),
make\_pair(``z'', 26), make\_pair(``e'', 5) \}; // default using
\textless{} operator auto minInt = min\_element(intVec.begin(),
intVec.end()); //Using pairLessThanFunction auto minPairFunction =
min\_element(pairVector.begin(), pairVector.end(),
pairLessThanFunction); //print minimum of intVector cout
\textless\textless{} ``min int from default:'' \textless\textless{}
\emph{minInt \textless\textless{} endl;  //print minimum of pairVector
cout \textless\textless{} ``min pair from PairLessThanFunction:''
\textless\textless{} (}minPairFunction).second \textless\textless{}
endl; return 0; \} Output min int from default: 6 min pair from
PairLessThanFunction: 2 Section 62.6: std::nd\_if template
\textless class InputIterator, class UnaryPredicate\textgreater{}
InputIterator find\_if (InputIterator first, InputIterator last,
UnaryPredicate pred); Eects Finds the rst element in a range for which
the predicate function pred returns true. Parameters first
=\textgreater{} iterator pointing to the beginning of the range last
=\textgreater{} iterator pointing to the end of the range pred
=\textgreater{} predicate function(returns true or false) Return An
iterator that points to the rst element within the range the predicate
function pred returns true for. The iterator points to last if val is
not found Example \#include \#include \#include using namespace std;
/\emph{ define some functions to use as predicates }/ //Returns true if
x is multiple of 10 bool multOf10(int x) \{ return x \% 10 == 0; \}
//returns true if item greater than passed in parameter class Greater \{
int \_than; public: Greater(int th):\_than(th)\{ \} bool operator()(int
data) const  \{ return data \textgreater{} \_than; \} \}; int main() \{
vector myvec \{2, 5, 6, 10, 56, 7, 48, 89, 850, 7, 456\}; //with a
lambda function vector::iterator gt10 = find\_if(myvec.begin(),
myvec.end(), \href{int\%20x}{}\{return x\textgreater10;\}); //
\textgreater= C++11 //with a function pointer vector::iterator pow10 =
find\_if(myvec.begin(), myvec.end(), multOf10); //with functor
vector::iterator gt5 = find\_if(myvec.begin(), myvec.end(), Greater(5));
//not Found vector::iterator nf = find\_if(myvec.begin(), myvec.end(),
Greater(1000)); // nf points to myvec.end() //check if pointer points to
myvec.end() if(nf != myvec.end()) \{ cout \textless\textless{} ``nf
points to:'' \textless\textless{} \emph{nf \textless\textless{} endl; \}
else \{ cout \textless\textless{} ``item not found''
\textless\textless{} endl; \} cout \textless\textless{} ``First item
\textgreater{} 10:'' \textless\textless{} }gt10 \textless\textless{}
endl; cout \textless\textless{} ``First Item n * 10:''
\textless\textless{} \emph{pow10 \textless\textless{} endl; cout
\textless\textless{} ``First Item \textgreater{} 5:''
\textless\textless{} }gt5 \textless\textless{} endl; return 0; \} Output
item not found First item \textgreater{} 10: 56 First Item n * 10: 10
First Item \textgreater{} 5: 6 Section 62.7: Using std::nth\_element To
Find The Median (Or Other Quantiles) The std::nth\_element algorithm
takes three iterators: an iterator to the beginning, nth position, and
end. Once the function returns, the nth element (by order) will be the
nth smallest element. (The function has more elaborate overloads, e.g.,
some taking comparison functors; see the above link for all the
variations.) Note This function is very ecient - it has linear
complexity. For the sake of this example, let's dene the median of a
sequence of length n as the element that would be in position n / 2. For
example, the median of a sequence of length 5 is the 3rd smallest
element, and so is the median of a sequence of length 6. To use this
function to nd the median, we can use the following. Say we start with
std::vector v\{5, 1, 2, 3, 4\};\\
std::vector::iterator b = v.begin(); std::vector::iterator e = v.end();
std::vector::iterator med = b; std::advance(med, v.size() / 2); // This
makes the 2nd position hold the median. std::nth\_element(b, med, e);\\
// The median is now at v{[}2{]}. To nd the pth quantile, we would
change some of the lines above: const std::size\_t pos = p *
std::distance(b, e); std::advance(nth, pos); and look for the quantile
at position pos. Section 62.8: std::count template \textless class
InputIterator, class T\textgreater{} typename
iterator\_traits::difference\_type count (InputIterator first,
InputIterator last, const T\& val); Eects Counts the number of elements
that are equal to val Parameters first =\textgreater{} iterator pointing
to the beginning of the range last =\textgreater{} iterator pointing to
the end of the range val =\textgreater{} The occurrence of this value in
the range will be counted Return The number of elements in the range
that are equal(==) to val. Example \#include \#include \#include using
namespace std; int main(int argc, const char * argv{[}{]}) \{  //create
vector vector intVec\{4,6,8,9,10,30,55,100,45,2,4,7,9,43,48\}; //count
occurrences of 9, 55, and 101 size\_t count\_9 = count(intVec.begin(),
intVec.end(), 9); //occurs twice size\_t count\_55 =
count(intVec.begin(), intVec.end(), 55); //occurs once size\_t
count\_101 = count(intVec.begin(), intVec.end(), 101); //occurs once
//print result cout \textless\textless{} ``There are''
\textless\textless{} count\_9 \textless\textless{} "
9s``\textless\textless{} endl; cout \textless\textless{}''There is "
\textless\textless{} count\_55 \textless\textless{} "
55``\textless\textless{} endl; cout \textless\textless{}''There is "
\textless\textless{} count\_101 \textless\textless{} "
101"\textless\textless{} ends; //find the first element == 4 in the
vector vector::iterator itr\_4 = find(intVec.begin(), intVec.end(), 4);
//count its occurrences in the vector starting from the first one
size\_t count\_4 = count(itr\_4, intVec.end(), \emph{itr\_4); // should
be 2 cout \textless\textless{} ``There are'' \textless\textless{}
count\_4 \textless\textless{} " " \textless\textless{} }itr\_4
\textless\textless{} endl; return 0; \} Output There are 2 9s There is 1
55 There is 0 101 There are 2 4 Section 62.9: std::count\_if template
\textless class InputIterator, class UnaryPredicate\textgreater{}
typename iterator\_traits::difference\_type count\_if (InputIterator
first, InputIterator last, UnaryPredicate red); Eects Counts the number
of elements in a range for which a specied predicate function is true
Parameters first =\textgreater{} iterator pointing to the beginning of
the range last =\textgreater{} iterator pointing to the end of the range
red =\textgreater{} predicate function(returns true or false) Return The
number of elements within the specied range for which the predicate
function returned true. Example \#include \#include \#include using
namespace std; /\emph{ Define a few functions to use as predicates }/
//return true if number is odd bool isOdd(int i)\{ return i\%2 == 1; \}
//functor that returns true if number is greater than the value of the
constructor parameter provided class Greater \{ int \_than; public:
Greater(int th): \_than(th)\{\} bool operator()(int i)\{ return i
\textgreater{} \_than; \} \}; int main(int argc, const char *
argv{[}{]}) \{ //create a vector vector myvec =
\{1,5,8,0,7,6,4,5,2,1,5,0,6,9,7\}; //using a lambda function to count
even numbers size\_t evenCount = count\_if(myvec.begin(), myvec.end(),
\href{int\%20i}{}\{return i \% 2 == 0;\}); // \textgreater= C++11
//using function pointer to count odd number in the first half of the
vector size\_t oddCount = count\_if(myvec.begin(), myvec.end()-
myvec.size()/2, isOdd); //using a functor to count numbers greater than
5 size\_t greaterCount = count\_if(myvec.begin(), myvec.end(),
Greater(5)); cout \textless\textless{} ``vector size:''
\textless\textless{} myvec.size() \textless\textless{} endl; cout
\textless\textless{} ``even numbers:'' \textless\textless{} evenCount
\textless\textless{} " found" \textless\textless{} endl; cout
\textless\textless{} ``odd numbers:'' \textless\textless{} oddCount
\textless\textless{} " found" \textless\textless{} endl; cout
\textless\textless{} ``numbers \textgreater{} 5:'' \textless\textless{}
greaterCount \textless\textless{} " found"\textless\textless{} endl;
return 0; \} Output vector size: 15 even numbers: 7 found odd numbers: 4
found numbers \textgreater{} 5: 6 found

\hypertarget{professional-notes-chapter-67-sorting-1}{%
\chapter{Professional Notes: Chapter 67:
Sorting}\label{professional-notes-chapter-67-sorting-1}}

Section 67.1: Sorting and sequence containers std::sort, found in the
standard library header algorithm, is a standard library algorithm for
sorting a range of values, dened by a pair of iterators. std::sort takes
as the last parameter a functor used to compare two values; this is how
it determines the order. Note that std::sort is not stable. The
comparison function must impose a Strict, Weak Ordering on the elements.
A simple less-than (or greater-than) comparison will suce. A container
with random-access iterators can be sorted using the std::sort
algorithm: Version C++11 \#include \#include std::vector MyVector = \{3,
1, 2\} //Default comparison of \textless{} std::sort(MyVector.begin(),
MyVector.end()); std::sort requires that its iterators are random access
iterators. The sequence containers std::list and std::forward\_list
(requiring C++11) do not provide random access iterators, so they cannot
be used with std::sort. However, they do have sort member functions
which implement a sorting algorithm that works with their own iterator
types. Version C++11 \#include \#include std::list MyList = \{3, 1, 2\}
//Default comparison of \textless{} //Whole list only. MyList.sort();
Their member sort functions always sort the entire list, so they cannot
sort a sub-range of elements. However, since list and forward\_list have
fast splicing operations, you could extract the elements to be sorted
from the list, sort them, then stu them back where they were quite
eciently like this: void sort\_sublist(std::list\& mylist,
std::list::const\_iterator start, std::list::const\_iterator end) \{
//extract and sort half-open sub range denoted by start and end iterator
std::list tmp; tmp.splice(tmp.begin(), list, start, end); tmp.sort();
//re-insert range at the point we extracted it from list.splice(end,
tmp); \} Section 67.2: sorting with std::map (ascending and descending)
This example sorts elements in ascending order of a key using a map. You
can use any type, including class, instead of std::string, in the
example below. \#include \#include \#include int main() \{
std::map\textless double, std::string\textgreater{} sorted\_map; // Sort
the names of the planets according to their size
sorted\_map.insert(std::make\_pair(0.3829, ``Mercury''));
sorted\_map.insert(std::make\_pair(0.9499, ``Venus''));
sorted\_map.insert(std::make\_pair(1, ``Earth''));
sorted\_map.insert(std::make\_pair(0.532, ``Mars''));
sorted\_map.insert(std::make\_pair(10.97, ``Jupiter''));
sorted\_map.insert(std::make\_pair(9.14, ``Saturn''));
sorted\_map.insert(std::make\_pair(3.981, ``Uranus''));
sorted\_map.insert(std::make\_pair(3.865, ``Neptune'')); for (auto
const\& entry: sorted\_map) \{ std::cout \textless\textless{}
entry.second \textless\textless{} " (" \textless\textless{} entry.first
\textless\textless{} " of Earth's radius)" \textless\textless{} `\n'; \}
\} Output: Mercury (0.3829 of Earth's radius) Mars (0.532 of Earth's
radius) Venus (0.9499 of Earth's radius) Earth (1 of Earth's radius)
Neptune (3.865 of Earth's radius) Uranus (3.981 of Earth's radius)
Saturn (9.14 of Earth's radius) Jupiter (10.97 of Earth's radius) If
entries with equal keys are possible, use multimap instead of map (like
in the following example). To sort elements in descending manner,
declare the map with a proper comparison functor
(std::greater\textless\textgreater): \#include \#include \#include int
main() \{ std::multimap\textless int, std::string,
std::greater\textgreater{} sorted\_map; // Sort the names of animals in
descending order of the number of legs
sorted\_map.insert(std::make\_pair(6, ``bug''));
sorted\_map.insert(std::make\_pair(4, ``cat''));
sorted\_map.insert(std::make\_pair(100, ``centipede''));
sorted\_map.insert(std::make\_pair(2, ``chicken''));
sorted\_map.insert(std::make\_pair(0, ``fish''));
sorted\_map.insert(std::make\_pair(4, ``horse''));
sorted\_map.insert(std::make\_pair(8, ``spider'')); for (auto const\&
entry: sorted\_map) \{ std::cout \textless\textless{} entry.second
\textless\textless{} " (has " \textless\textless{} entry.first
\textless\textless{} " legs)" \textless\textless{} `\n';  \} \} Output
centipede (has 100 legs) spider (has 8 legs) bug (has 6 legs) cat (has 4
legs) horse (has 4 legs) chicken (has 2 legs) fish (has 0 legs) Section
67.3: Sorting sequence containers by overloaded less operator If no
ordering function is passed, std::sort will order the elements by
calling operator\textless{} on pairs of elements, which must return a
type contextually convertible to bool (or just bool). Basic types
(integers, oats, pointers etc) have already build in comparison
operators. We can overload this operator to make the default sort call
work on user-dened types. // Include sequence containers \#include
\#include \#include // Insert sorting algorithm \#include \\
class Base \{ public: // Constructor that set variable to the value of v
Base(int v): variable(v) \{ \} // Use variable to provide total order
operator less //\passthrough{\lstinline!this!} always represents the
left-hand side of the compare. bool operator\textless(const Base \&b)
const \{ return this-\textgreater variable \textless{} b.variable; \}
int variable; \}; int main() \{ std::vector vector; std::deque deque;
std::list list; // Create 2 elements to sort Base a(10); Base b(5); //
Insert them into backs of containers vector.push\_back(a);
vector.push\_back(b);  deque.push\_back(a); deque.push\_back(b);
list.push\_back(a); list.push\_back(b); // Now sort data using
operator\textless(const Base \&b) function std::sort(vector.begin(),
vector.end()); std::sort(deque.begin(), deque.end()); // List must be
sorted differently due to its design list.sort(); return 0; \} Section
67.4: Sorting sequence containers using compare function // Include
sequence containers \#include \#include \#include // Insert sorting
algorithm \#include class Base \{ public: // Constructor that set
variable to the value of v Base(int v): variable(v) \{ \} int variable;
\}; bool compare(const Base \&a, const Base \&b) \{ return a.variable
\textless{} b.variable; \} int main() \{ std::vector vector; std::deque
deque; std::list list; // Create 2 elements to sort Base a(10); Base
b(5); // Insert them into backs of containers vector.push\_back(a);
vector.push\_back(b); deque.push\_back(a); deque.push\_back(b);
list.push\_back(a); list.push\_back(b); // Now sort data using comparing
function  std::sort(vector.begin(), vector.end(), compare);
std::sort(deque.begin(), deque.end(), compare); list.sort(compare);
return 0; \} Section 67.5: Sorting sequence containers using lambda
expressions (C++11) Version C++11 // Include sequence containers
\#include \#include \#include \#include \#include // Include sorting
algorithm \#include class Base \{ public: // Constructor that set
variable to the value of v Base(int v): variable(v) \{ \} int variable;
\}; int main() \{ // Create 2 elements to sort Base a(10); Base b(5); //
We're using C++11, so let's use initializer lists to insert items.
std::vector vector = \{a, b\}; std::deque deque = \{a, b\}; std::list
list = \{a, b\}; std::array \textless Base, 2\textgreater{} array = \{a,
b\}; std::forward\_list flist = \{a, b\}; // We can sort data using an
inline lambda expression std::sort(std::begin(vector), std::end(vector),
\href{const\%20Base\%20\&a,\%20const\%20Base\%20\&b}{} \{ return
a.variable \textless{} b.variable;\}); // We can also pass a lambda
object as the comparator // and reuse the lambda multiple times auto
compare = \href{const\%20Base\%20\&a,\%20const\%20Base\%20\&b}{} \{
return a.variable \textless{} b.variable;\};
std::sort(std::begin(deque), std::end(deque), compare);
std::sort(std::begin(array), std::end(array), compare);
list.sort(compare); flist.sort(compare); return 0; \} Section 67.6:
Sorting built-in arrays The sort algorithm sorts a sequence dened by two
iterators. This is enough to sort a built-in (also known as c- style)
array. Version C++11 int arr1{[}{]} = \{36, 24, 42, 60, 59\}; // sort
numbers in ascending order sort(std::begin(arr1), std::end(arr1)); //
sort numbers in descending order sort(std::begin(arr1), std::end(arr1),
std::greater()); Prior to C++11, end of array had to be ``calculated''
using the size of the array: Version \textless{} C++11 // Use a
hard-coded number for array size sort(arr1, arr1 + 5); // Alternatively,
use an expression const size\_t arr1\_size = sizeof(arr1) /
sizeof(*arr1); sort(arr1, arr1 + arr1\_size); Section 67.7: Sorting
sequence containers with specifed ordering If the values in a container
have certain operators already overloaded, std::sort can be used with
specialized functors to sort in either ascending or descending order:
Version C++11 \#include \#include \#include std::vector v =
\{5,1,2,4,3\}; //sort in ascending order (1,2,3,4,5)
std::sort(v.begin(), v.end(), std::less()); // Or just:
std::sort(v.begin(), v.end()); //sort in descending order (5,4,3,2,1)
std::sort(v.begin(), v.end(), std::greater()); //Or just:
std::sort(v.rbegin(), v.rend()); Version C++14 In C++14, we don't need
to provide the template argument for the comparison function objects and
instead let the object deduce based on what it gets passed in:
std::sort(v.begin(), v.end(), std::less\textless\textgreater()); //
ascending order std::sort(v.begin(), v.end(),
std::greater\textless\textgreater()); // descending order

\begin{center}\rule{0.5\linewidth}{0.5pt}\end{center}
